% \documentclass[a4paper,12pt]{article}{{{
\documentclass[a4paper,11pt]{article}
% -------------------------------------------------
% Set up the page margins.
% -------------------------------------------------
% \usepackage[text={6.5in,9.5in},centering]{geometry}
% \setlength{\topmargin}{-0.25in}
\usepackage[margin=1in]{geometry}

% -------------------------------------------------
% Load Packages here
% -------------------------------------------------
%\usepackage{hyperref}
\usepackage{graphicx,url,subfig}
\RequirePackage[OT1]{fontenc}
\RequirePackage{amsthm,amsmath,amssymb,amscd}
\RequirePackage[numbers]{natbib}
\RequirePackage[colorlinks,citecolor=blue,urlcolor=blue]{hyperref}
\usepackage{graphics}
\usepackage{enumerate}
\usepackage{datetime}
\usepackage{etex}
% \usepackage[notcite,notref,color]{showkeys}
\usepackage[normalem]{ulem} % either use this (simple) or
\usepackage{soul} % use this (many fancier options)
\makeatother
\numberwithin{equation}{section}
\allowdisplaybreaks
\usepackage{mathrsfs}

%copied from invariant measure paper
%-----
\usepackage{amssymb, latexsym, amsmath, graphics, fullpage, epsfig, amsthm, relsize, pgf, tikz, amsfonts, makeidx, latexsym, ifthen, calc}
\usepackage[colorlinks,citecolor=blue,urlcolor=blue]{hyperref}
%\usepackage{eucal} this is a different font than \mathcal{•}
\usetikzlibrary{arrows}
\usepackage{mathrsfs}
\usepackage[shortlabels]{enumitem}
\usepackage{tikz}
\usepackage{amsmath}
\usetikzlibrary{arrows}
\usetikzlibrary{shadows.blur}
\usepackage{pgfplots}
\usepackage{mwe} % For dummy images
% \usepackage{subcaption}
\usepackage{multirow}
\usepackage{dcolumn}
\newcolumntype{2}{D{.}{}{2.0}}
\usepackage{multicol}
\numberwithin{equation}{section}
%----
%end copy from invariant measure paper

% -------------------------------------------------
% check references
% -------------------------------------------------
% \usepackage{refcheck}
% \usepackage[notref,notcite]{showkeys}
% \usepackage[notcite]{showkeys}
% \usepackage{showkeys}

% -------------------------------------------------
% Environments here
% -------------------------------------------------
\theoremstyle{plain}
\newtheorem{theorem}{Theorem}[section]
\newtheorem{corollary}[theorem]{Corollary}
\newtheorem{lemma}[theorem]{Lemma}
\newtheorem{proposition}[theorem]{Proposition}
\newtheorem{exercise}{Exercise}
\theoremstyle{definition}
\newtheorem{definition}[theorem]{Definition}
\newtheorem{hypothesis}[theorem]{Hypothesis}
\newtheorem{assumption}[theorem]{Assumption}

\newtheorem{remark}[theorem]{Remark}
\newtheorem{conjecture}[theorem]{Conjecture}
\newtheorem{example}[theorem]{Example}

% -------------------------------------------------
%  Define short hand symbols.
% -------------------------------------------------
\newcommand{\E}{\mathbb{E}}
\newcommand{\W}{\dot{W}}
\newcommand{\B}{\textbf{B}}
\newcommand{\D}{\mathbb{D}}
\newcommand{\ud}{\ensuremath{\mathrm{d} }}
\newcommand{\Ceil}[1]{\left\lceil #1 \right\rceil}
\newcommand{\Floor}[1]{\left\lfloor #1 \right\rfloor}
\newcommand{\sgn}{\text{sgn}}
\newcommand{\Lad}{\text{L}_{\text{ad}}^2}
\newcommand{\SI}[1]{\mathcal{I}\left[#1 \right]}
\newcommand{\SIB}[2]{\mathcal{I}_{#2}\left[#1 \right]}
\newcommand{\Indt}[1]{1_{\left\{#1 \right\} }}
\newcommand{\LadInPrd}[1]{\left\langle #1 \right\rangle_{\text{L}_\text{ad}^2}}
\newcommand{\LadNorm}[1]{\left|\left|  #1 \right|\right|_{\text{L}_\text{ad}^2}}
\newcommand{\Norm}[1]{\left\|  #1   \right\|}
\newcommand{\Ito}{It\^{o} }
\newcommand{\Itos}{It\^{o}'s }
\newcommand{\spt}[1]{\text{supp}\left(#1\right)}
\newcommand{\InPrd}[1]{\left\langle #1 \right\rangle}
\newcommand{\mr}{\textbf{r}}
\newcommand{\Ei}{\text{Ei}}
\newcommand{\arctanh}{\operatorname{arctanh}}
\newcommand{\ind}[1]{\mathbb{I}_{\left\{ {#1} \right\} }}

\DeclareMathOperator{\esssup}{\ensuremath{ess\,sup}}

\newcommand{\steps}[1]{\vskip 0.3cm \textbf{#1}}
\newcommand{\calB}{\mathcal{B}}
\newcommand{\calD}{\mathcal{D}}
\newcommand{\calE}{\mathcal{E}}
\newcommand{\calF}{\mathcal{F}}
\newcommand{\calG}{\mathcal{G}}
\newcommand{\calK}{\mathcal{K}}
\newcommand{\calH}{\mathcal{H}}
\newcommand{\calI}{\mathcal{I}}
\newcommand{\calL}{\mathcal{L}}
\newcommand{\calM}{\mathcal{M}}
\newcommand{\calN}{\mathcal{N}}
\newcommand{\calO}{\mathcal{O}}
\newcommand{\calT}{\mathcal{T}}
\newcommand{\calP}{\mathcal{P}}
\newcommand{\calR}{\mathcal{R}}
\newcommand{\calS}{\mathcal{S}}
\newcommand{\bbC}{\mathbb{C}}
\newcommand{\bbN}{\mathbb{N}}
\newcommand{\bbP}{\mathbb{P}}
\newcommand{\bbZ}{\mathbb{Z}}
\newcommand{\myVec}[1]{\overrightarrow{#1}}
\newcommand{\sincos}{\begin{array}{c} \cos \\ \sin \end{array}\!\!}
\newcommand{\CvBc}[1]{\left\{\:#1\:\right\}}
\newcommand*{\one}{ {{\rm 1\mkern-1.5mu}\!{\rm I} }}
\newcommand{\RR}{\mathbb{R}}
\newcommand{\R}{\mathbb{R}}
\newcommand{\myEnd}{\hfill$\square$}
\newcommand{\ds}{\displaystyle}
\newcommand{\Shi}{\text{Shi}}
\newcommand{\Chi}{\text{Chi}}
\newcommand{\Daw}{\ensuremath{\mathrm{daw} }}
\newcommand{\Erf}{\ensuremath{\mathrm{erf} }}
\newcommand{\Erfi}{\ensuremath{\mathrm{erfi} }}
\newcommand{\Erfc}{\ensuremath{\mathrm{erfc} }}
\newcommand{\He}{\ensuremath{\mathrm{He} }}

\def\crtL{\theta}

\newcommand{\conRef}[1]{(#1)}

\DeclareMathOperator{\Lip}{\mathit{L}}
\DeclareMathOperator{\LIP}{Lip}
\DeclareMathOperator{\lip}{\mathit{l}}

\DeclareMathOperator{\Vip}{\overline{\varsigma}}
\DeclareMathOperator{\vip}{\underline{\varsigma}}
\DeclareMathOperator{\vv}{\varsigma}
% }}}

\begin{document}

\title{Fractional SPDE with Time-independent Noise:\\
Solvability and Exact Asymptotics\footnote{Department of Mathematics and
Statistics, Auburn University, Auburn, Alabama.}}
\author{Le Chen\footnote{Email: \url{lzc0090@auburn.edu}} and Nick Eisenberg\footnote{Corresponding author. Email:
\url{...@auburn.edu}}}
\date{\today}

\maketitle

\begin{abstract}
	\noindent In this article, we study the ......
\end{abstract}

\noindent {\em MSC 2010:} Primary 60H15; Secondary 60H07, 37H15

%60H15=SPDEs
%60H07 stochastic calculus of variations and the Malliavin calculus
%37H15 Multiplicative ergodic theory, Lyapunov exponents
\vspace{1mm}

\noindent {\em Keywords:}
Stochastic partial differential equations; Caputo derivatives;
Riemann-Liouville fractional integral; the Fox H-function; Malliavin
calculus; Lyapunov exponents; Exact moment asymptotics; Moment blowup.

\tableofcontents
% \pagebreak
\section{Introduction}
In this paper we study the following fractional stochastic equation:
	\begin{equation} \label{spdee}
		\begin{cases}
			\left(\partial^b + \frac{\nu}{2}(-\Delta)^{\alpha/2}\right)u(t,x) = I^\gamma_t \left[\sqrt{\theta}u(t,x)\dot W(x) \right] & \text{$x\in \R^d$ , $t>0$} \\
			u(0,\cdot) = 1                                                                                                         & b \in (0,1]            \\
			u(0,\cdot) = 1, \quad u_t(0,\cdot) = 		0                                                                                 & b \in (1,2]
		\end{cases}
	\end{equation}
where $a \in (0,2)$, $b \in (0,2)$, $\nu > 0$, $r>0$ and $\theta > 0$ and $W = \left\{ W(\varphi) : \varphi \in \mathcal{D}(\R^d)  \right\}$ is a time-independent Gaussian process, defined on a complete probability space $(\Omega, \mathcal{F}, P)$, with mean zero and covariance
	\[
		\mathbb{E}[ W(\varphi) W(\psi) ] = \int_{\R^d} \mathcal{F}\varphi(\xi)\overline{\mathcal{F}\psi(\xi)} \mu(\ud \xi) =: \langle \varphi, \psi \rangle_{\mathcal{H}},
	\]
with the spectral measure $\mu$ being assumed to be a non-negative an non-negative definite tempered measure on $\R^d$. We choose $\mathcal{F}\varphi(\xi) = \int_{\R^d} \exp(-ix \xi) \varphi(x) \ud x$ to be the Fourier transform of $\varphi$. It is known that the Fourier transform of $\mu$, denoted by $\gamma$, is also a non-negative and non-negative definite tempered measure. See pages 51 and 52 of \cite{DalangEtc09Minicourse}.

Above, $\partial^b_t$ denotes the Caputo fractional differential operator:
	\[ \partial^b_tf(t) :=
		\begin{cases}
			\dfrac{1}{\Gamma(m-b)} \int_0^t \ud \tau \dfrac{f^{m}(\tau)}{(t-\tau)^{b+1-m}} & \text{if } m-1 < b < m,
			\\\dfrac{\ud^m}{\ud t^m}f(t) & \text{if } b=m,
		\end{cases}
	\]
and $I^r_t$ is the Riemann-Liouville fractional integral of order $r>0$:
	\[
		I^r_tf(t) := \frac{1}{\Gamma(r)} \int_0^t (t-s)^{r-1}f(s)\ud s, \quad \text{for } t>0,
	\]
with the convention that $I^0_t = \text{Td (the identity operator)}$.

 The fundamental solutions to \ref{spdee} consist of a triplet:
 	\[
 		\left\{ Z_{a,b,\nu,d}(t,x),  Z^*_{a,b,\nu,d}(t,x),  Y_{a,b,r,\nu,d}(t,x) : t >0, x \in \R^d \right\}.
 	\]
Each member of the triplet is based on the Fox H-function which is much more technically complicated than the Green's function for either the heat or wave equation. The actual meanings of these are not to important in this setting so for that reason we direct the reader to \cite{ChenHuNualartFraction} for the definitions. Since we are interested in the constant one initial condition (and 0 initial velocity when $b>1$), by Theorem 4.1 \cite{ChenHuNualartFraction} the corresponding solution to the homogeneous equation (i.e. the solution when there is no driving source) is equal to the constant one. Hence, by setting $G(t,x) =  Y_{a,b,r,\nu,d}(t,x)$, \ref{spdee} can be written as the following Skorohod integral equation:
	\begin{equation}\label{spdeint}
		u(t,x) = 1 + \sqrt{\theta}\int_0^t \left( \int_{\R^d}G(t-s,x-y)u(s,y) W(\delta y) \right) \ud s.
	\end{equation}

We now need to make mention of the suitable choices for the noise $W$. Let $k \in \{1, \cdots , d\}$ and partition the $d$-coordinates of $x=(x_1, \cdots, x_n)$ into $k$ distinct groups of size $d_i$ for $i \in \{1, \cdots, k\}$. Note that $d_1 + \cdots + d_k = d$. Denote $x_{(i)} = (x_{i_1} , \cdots , x_{i_{d_i}})$ to be the coordinates in the $i^{th}$ partition. We consider a Gaussian Noise with correlation function  given by 
	\begin{equation}\label{corfunction}
		\gamma(x) = \prod_{i=1}^k |x_{(i)}|^{-\alpha_i}
	\end{equation}
where $\alpha_i \in (0, d_i)$ and we define $\alpha = \sum_{i=1}^k \alpha_i$. This correlation function corresponds with the spectral density given by $\mu (\ud \xi) = \varphi(\xi) \ud \xi$ with 
	\begin{equation}\label{specfunction}
		\varphi(\xi) = C\prod_{i=1}^k|\xi_{(i)}|^{-(d_i-\alpha_i)}.
	\end{equation}
The Constant $C$ can be calculated through the relation $\mathcal{F}\gamma = \varphi$. An important fact stated in \cite{BalanChenChen} is that there exists a non-negative function $K$ on $\R^d$ such that $\gamma = K*K$ where '$*$' denotes the spatial convolution.

	\begin{example} 
When $k=1$ above then we recover the Riesz kernel $\gamma(x) = |x|^{-\alpha}$ and its spectral density $\varphi(x) = C_{d,\alpha}|x|^{-(d-\alpha)}$. In addition, we may write $\gamma = K*K$ where $"*"$ denotes the convolution and $K(x) = \beta_{\alpha,d}|x|^{-(d+\alpha)/2}$. In this case we have 		
		\begin{equation}\label{noiseconstant}
			C_{d,\alpha} = \pi^{-d/2}2^{-\alpha} \frac{\Gamma((d-\alpha)/2)}{\Gamma(\alpha/2)} \quad    \text{and}	\quad	\beta_{\alpha,d} = \pi^{-d/4}\frac{\Gamma((d+\alpha)/4)}{\Gamma((d-\alpha)/4)}\sqrt{ \frac{\Gamma((d-\alpha)/2)}{\Gamma(\alpha/2)} }.
		\end{equation}
	\end{example}

	\begin{example}
When $k=d$ above then we recover the noise, $W$, with correlation function given by $\gamma(x) = \prod_{i=1}^d |x_i|^{-\alpha_i}$ where $\alpha_i \in (0,1)$. In this case, the spectral density is given by $\varphi(\xi) = \prod_{i=1} C_{1,\alpha_i} |\xi_i|^{-(1-\alpha_i)}$. We may also write $\gamma = K*K$ where $K(x)=\prod_{i=1}^d \beta_{\alpha_i,1}|x_i|^{-(1+\alpha_i)/2}$ where $C_{1,\alpha_i}$ and $ \beta_{\alpha_i,1}$ are given by \ref{noiseconstant}. If we consider the parameterization $\alpha_i = 2-2H_i$ with $H_i \in (1/2 , 1)$, then the corresponding noise is the fractional Brownian sheet with indicies $H_1, \cdots , H_d$. The heat equation with this noise was studied in \cite{HuFracBrown}.
	\end{example}

	\begin{remark}\label{spectralconjecture}
We may actually apply the techniques in this paper to any such noise whose correlation function and spectral density, $\gamma$ and $\mu$, satisfy
		\begin{enumerate}
			\item both $\mu$ and $\gamma$ are absolutely continuous with respect to Lebesgue measure
			\item for some $\alpha \in (0,d)$, $\gamma$ satisfies that for all $c>0$ and $x\in \R^d$ that $\gamma(cx) = c^{-\alpha}\gamma(x)$
			\item there exists a nonnegative function $K$ on $\R^d$ such that $\gamma = K*K$ where $"*"$ refers to the spatial convolution. 
		\end{enumerate}
However, there is an issue with applying the results in this paper to the more general noises satisfying \ref{spectralconjecture}. Proving $\rho$, defined below, is finite is a challange. In \cite{chenlarge}, $\rho$ is proven to be finite in the case where $\varphi$ is the Riesz kernel. Below in Lemma \ref{rhofinite}, we give a proof showing that $\rho$ is finite for the noise with spectral density defined above in \ref{specfunction}.

\textcolor{red}{Are there other noises satisfying \ref{spectralconjecture} other than the general fractional noise?}
	\end{remark}

Next, we need to make some assumption on the parameters $a,b,r$ of \ref{spdee} as well as the dimension $d$. Under theses assumptions, the Greens function $G$ will be non-negative and smooth which was shown in Theorem 4.6 \cite{ChenHuNualartFraction} and equation (4.5) \cite{ChenHuNualartFraction} respectively.
	
	\begin{assumption}\label{paramassumptions}
The parameters $a,b,r$ and the dimension $d$ must satisfy one of the following conditions:
		\begin{enumerate}
			\item $d \ge 1$, $b\in(0,1)$, $a\in(0,2]$, $r\ge0$;
			\item $d \ge 1$, $b=1$, $a\in (0,2]$, $r=0$ or $r>1$;
			\item $1\le d \le 3$, $1<b<a\le 2$, $r>0$;
			\item $1 \le d \le 3$, $1 < b=a < 2$, $r > \dfrac{d+3}{2}-b$.
		\end{enumerate}
	\end{assumption}


We now introduce the operator 
	\[
		\mathcal{M}_a(\Theta f , \theta) = \sup_{g \in \mathcal{F}_a} \left\{ \langle g^2 * g^2, f \rangle^{1/2} - \frac{\theta}{2}\mathcal{E}_a(g,g) \right\}
	\]
and define $\mathcal{M}_a := \mathcal{M}_a(|\cdot|^{-\alpha},1)$. See Subsection \ref{largedevresults} of this paper for more details. Our goal is to develop exact moment asymptotics to the solution $u(t,x)$. We follow the techniques laid out in \cite{BalanChenChen} where in that paper they focus exclusively on the stochastic wave equation. This paper acts as a generalization of that paper to the fraction stochastic equation. Here is the main result of this paper.

	\begin{theorem}\label{mainresult}
Under Assumption \ref{paramassumptions}, equation \ref{spdee} driven by the noise with spectral density given in equation \ref{specfunction} has a unique solution $u(t,x)$ in $L^p(\Omega)$ for any $p \ge 2$, $t>0$ and $x \in \R^d$ if 
		\begin{equation}
			\alpha < \min\left( \frac{a}{b}\left[2(b+r)-1\right], 2a, d \right).
		\end{equation} 
Moreover, we have the following moment asymptotic: 
		\begin{equation} \label{momentasym}
			\lim_{t_p \to \infty} t_p^{-\beta} \log \Norm{u(t,x)}_p = \left(\frac{1}{2}\right)\left( \frac{2a}{2a(b+r)- b\alpha} \right)^\beta (\theta\nu^{-\alpha/a} \mathcal{M}_a^{\frac{2a-\alpha}{a}})^{\frac{a}{2a(b+r)-b\alpha-a}}\left(2(b+r)-\frac{b\alpha}{a}-1\right).
		\end{equation}	
\textcolor{red}{I dont know the best way to organize and combine the like terms for the moment asymptotic \ref{momentasym}. Above is one way.}		
		
where
	\[
		\beta := \dfrac{2(b+r)-\dfrac{b\alpha}{a}}{ 2(b+r)-\dfrac{b\alpha}{a} -1 } \quad \text{and} \quad t_p^\beta := (p-1)^{\beta -1} t^\beta
	\]
	\end{theorem}
The existence and uniqueness of a $L^p(\Omega)$ solution will be proven in Section \ref{solvability}. The upper bound and lower bound for \ref{momentasym} will be proven in Sections \ref{upperbound} and \ref{lowerbound} respectively.
\section{Some preliminaries.}

\subsection{Malliavin Calculus}
To briefly summarize \cite{DalangEtc09Minicourse}, we start with a non-negative and non-negative definite tempered measure $\Gamma$ with density $\gamma$ in the sense that $\Gamma(\ud x) = \gamma(x) \ud x$, and 
	\[
		\int_{\R^d} \Gamma(\ud x) (\varphi * \tilde{\varphi})(x) \ge 0 \quad \text{for all $\varphi \in \mathscr{S}(\R^d)$}
	\]
where $\tilde{\varphi(x)} = \varphi(-x)$. In addition we require the existance of an $r>0$ such that 
	\[
		\int_{\R^d}\Gamma(\ud x) \frac{1}{ (1 + |x|^2 )^r} < \infty.
	\]
According to the Bochner-Schwarz theorem, there exists a non-negative measure $\mu$, often referred to as the spectral measure on $\R^d$ whose Fourier transform is $\Gamma$ in the sense that for any $\varphi \in \mathscr(\R^d)$ 
	\[
		\int_{\R^d}\Gamma(\ud x) \varphi(x) = \int_{\R^d} \mu(\ud t) \mathcal{F}\varphi (t) 
	\]

The functional defined on the space of test functions $\mathcal{D}(\R^d)$
	\[
		C(\varphi,\psi) = \int_{\R} \mathcal{F}\varphi(\xi) \overline{ \mathcal{F}\psi(\xi) } \mu(\ud \xi)
	\]
is non-negative define and thus we may associate it as the covariance functional of a zero-mean Gaussian processes, $W:= \{ W(\varphi) : \varphi \in \mathcal{D}(\R^d) \}$. In other words, 
	\[
		\mathbb{E}(W(\varphi)W(\psi)) = \int_{\R} \mathcal{F}\varphi(\xi) \overline{ \mathcal{F}\psi(\xi) } \mu(\ud \xi) =: \langle \varphi, \psi \rangle_{\mathcal{H}}.
	\]
Let $\mathcal{H}$ be the completion of $\mathcal{D}(\R^d)$ with respect to $\langle \cdot, \cdot \rangle_{\mathcal{H}}$ and thus we see $\varphi \mapsto W(\varphi)$ is an isometry from the normed space $(\mathcal{D}(\R^d), \Norm{\cdot}_{\mathcal{H}})$ to $L^2(\Omega)$ where $\Norm{f}_{\mathcal{H}} = \langle f, f \rangle_{\mathcal{H}}$.

We denote $\delta$ the Skorohod integral with respect to $W$ and we denote its domain $\text{Dom}(\delta)$. For any $u\in \text{Dom}(\delta)$, we write
	\[
		\delta(u) = \int_{\R^d} u(x) W(\delta x).
	\] 
Through construction of $\delta$, $\delta(u)$ is a random variable in $L^2(\Omega)$. For a deeper explanation, we direct the reader to Nualart \cite{nualartmalliavin}.

\begin{definition}
	A random field $u=\{u(t,x) : t \ge 0, x\in\R^d\}$ such that $\Norm{u(t,x)}_2 < \infty$ for all $t \ge 0$ and $ x \in \R^d$ is called a mild solution to the equation \ref{spdee} if the random field $\{G(t-s,x-y)u(s,y): s \in [0,t], y \in \R^d \}$ is Skorohod integrable for each $s\in [0,t]$ and $x\in \R^d$ fixed and satisfies the following integral equation:
		\begin{equation}
		u(t,x) = 1 + \sqrt{\theta}\int_0^t \left( \int_{\R^d} G(t-s,x-y)u(s,y) W(\delta y) \right)\ud s
		\end{equation}
\end{definition}

Through a Piccard Iteration Scheme we can see that the solution should be given by
	\begin{equation} \label{solwexp}
u(t,x) = 1 + \sum_{k=1}^\infty \theta^{k/2}I_k(f_k(\cdot,x,t))
	\end{equation}
where $I_k : \mathcal{H}^{\otimes k}((\R^d)^k) \to \mathcal{H}_k$ is the $k^{th}$ order Skorohod integral and $\mathcal{H}_k$ is the $k^{th}$ Wiener chaos space. The kernels $f_k(\cdot,x,t) \in \mathcal{H}^{\otimes k}((\R^d)^k)$ are calculated through the Piccard iteration scheme to be
\begin{align}\label{wker}
f_n(x_1,...,x_k;x,t) &= \int_0^t \int_0^{t_n}\cdots \int_0^{t_2} G(t-t_n,x-x_n)\cdots G(t_2-t_1,x_2-x_1) \ud t_1\cdots \ud t_n
\\&= \int_0^t \int_0^{t_n}\cdots \int_0^{t_2} G(t_1,x-x_n)\cdots G(t_n-t_{n-1},x_2-x_1) \ud t_1\cdots \ud t_n. \nonumber
\end{align}
For ease of notation throughout this article, we write the above integrals as
	\begin{align*}
		f_n(x_1,...,x_n;x,t) &= \int_{[0,t]^n_<} G(t-t_n,x-x_n)\cdots G(t_2-t_1,x_2-x_1) \ud \vec{t}
		\\&= \int_{[0,t]^n_<}G(t_1,x-x_n)\cdots G(t_n-t_{n-1},x_2-x_1) \ud \vec{t}.
	\end{align*}
where in general $\ud\vec{z} = \prod_{k=1}^d \ud z_k$ and the dimension will always be easily understood. We denote by $\tilde{f_n}(\cdot, x, t)$ to be the symmetrization of $f_n(\cdot,x,t)$ in the variables $x_1,\cdots,x_n$:
	\begin{align}
		\tilde{f_n}(\cdot; x, t) &= \frac{1}{n!} \sum_{\rho \in \Sigma_n} f_n(x_{\rho(1)},\cdots, x_{\rho(n)})
		\\&= \frac{1}{n!} \sum_{\rho \in \Sigma_n} \int_{[0,t]^n_<} G(t-t_n,x-x_{\rho(n)})\cdots G(t_2-t_1,x_{\rho(2)}-x_{\rho(1)}) \ud \vec{t}
	\end{align}
where $\Sigma_n$ is the set of all permutations of $\{1,\cdots,n\}  $. In particular, for $x=0$ we have
	\begin{align} \label{fperm0}
			\tilde{f_n}(\cdot; 0, t) = \frac{1}{n!} \sum_{\rho \in \Sigma_n} \int_{[0,t]^n_<} G(t_1,x_{\rho(1)})\cdots G(t_n-t_{n-1},x_{\rho(n)}-x_{\rho(n-1)}) \ud \vec{t}.
	\end{align}
The Fourier transform of the kernels, $f_n$, is given by
	\begin{align} \label{fweker}
		\mathcal{F}f_n(\cdot;x,t)(\xi_1,...,\xi_n) = e^{-i \left(\sum_{j=1}^n \xi_j \right) \cdot x} \int_{[0,t]^n <} \prod_{k=1}^n \overline{\mathcal{F}G(t_{t_k+1}-t_k,\cdot)\left(\sum_{j=1}^k \xi_j\right) } 	\ud \textbf{t}
	\end{align}
where $t_{n+1} = t$

Recall the notation above that $G(t,x) = Y_{a,b,r,d}(t,x)$. The following result shows an important scaling property that both $\mathcal{F}G(t,\cdot)(\xi)$ and $\Norm{\tilde{f_n}(\cdot,t,x)}_{\mathcal{H}^{\otimes n}}$ have.

	\begin{theorem}
		Let $c > 0$, then the two following scaling properties hold:
			\begin{equation}\label{greenscale}
			\mathcal{F}G(t,\cdot)(c\xi) = c^{-\frac{a}{b}(b+r -1)}\mathcal{F}G\left(c^{\frac{a}{b}}t,\cdot\right)( \xi)  \quad \text{or} \quad	\mathcal{F}G(ct,\cdot)(\xi) = c^{b+r -1}\mathcal{F}G(t,\cdot)(c^{b/a} \xi)
			\end{equation}
		and
			\begin{equation}\label{normscale}
					\Norm{\tilde{f_n}(\cdot,0,t)}_{H^{\otimes n}}^2 = t^{[2(b + r ) -b \alpha / a]n} \Norm{\tilde{f_n}(\cdot,0,1)}_{H^{\otimes n}}^2
			\end{equation}
	\end{theorem}
\noindent\textbf{Proof:}
The Fourier transform of $G$ is given in \cite{ChenHuNualartFraction}. With this we can directly calculate
	\begin{align*}
		   \mathcal{F}G(t, \cdot)(c\xi) & = t^{b + r -1} \sum_{k=0}^\infty \dfrac{\left(-2^{-1}\nu t^b |c\xi|^a\right)^k}{\Gamma\left(b k + (b + r)\right)} \\
		& = t^{b + r -1} \sum_{k=0}^\infty \dfrac{\left(-2^{-1}\nu t^b c^a |\xi|^a\right)^k}{\Gamma\left(b k + (b + r)\right)}\\
		& = t^{b + r -1} \sum_{k=0}^\infty \dfrac{\left(-2^{-1}\nu (tc^{\frac{a}{b}})^b |\xi|^\alpha\right)^k}{\Gamma\left(\beta k + (\beta + \gamma)\right)}\\
		& =c^{\frac{a}{b}(1-r-b)} (tc^{\frac{a}{b}})^{b + r -1} \sum_{k=0}^\infty \dfrac{\left(-2^{-1}\nu (tc^{\frac{a}{b}})^b |\xi|^a\right)^k}{\Gamma\left(b k + (b + \gamma)\right)}\\
		& =c^{\frac{a}{b}(1-r-b)} \mathcal{F}G\left(c^{\frac{a}{b}}t, \cdot\right)(\xi)
	\end{align*}
and this gives us the equality on \ref{greenscale}. The second equality in \ref{greenscale} is proven by letting $c=c^{b/a}$ in the first equality.

For \ref{normscale}, first recall that
	\begin{align*}
		\Norm{\tilde{f_n}(\cdot,0,t)}_{H^{\otimes n}}^2 =  \int_{\R^{nd}} |\mathcal{F}\tilde{f}_n(\cdot,0,t)|^2(\vec{\xi}) \varphi(\vec{\xi})\ud \vec{\xi}
	\end{align*}
where $\varphi(\vec{\xi}) = \varphi(\xi_1)\cdots\varphi(\xi_n)$. We start by applying this definition of $\tilde{f_n}$, \ref{fweker} and the scaling property \ref{normscale} to see
	\begin{align*}
		&  \int_{\R^{nd}} |\mathcal{F}\tilde{f}_n(\cdot,0,t)|^2(\vec{\xi}) \varphi(\vec{\xi})\ud \vec{\xi} =  \int_{\R^{nd}} \left| \frac{1}{n!} \sum_{\sigma \in \Sigma_n}\int_{[0,t]^n<} \overline{\mathcal{F}G(t-t_n,\cdot)(\sum_{i=1}^n\xi_i)\cdots \mathcal{F}G(t_2-t_1,\cdot)(\xi_1)} \right|^2 \varphi(\vec{\xi})\ud \vec{\xi}
		\\&= t^{2n(b+r-1)}\int_{\R^{nd}} \left|\frac{1}{n!} \sum_{\sigma \in \Sigma_n} \int_{[0,t]^n<} \overline{\mathcal{F}G\left(1-\frac{t_n}{t},\cdot\right)(t^{b/a}(\sum_{i=1}^n\xi_i))\cdots \mathcal{F}G\left(\frac{t_2}{t}-\frac{t_1}{t},\cdot\right)(t^{b/a}\xi_1)} \right|^2 \varphi(\vec{\xi})\ud \vec{\xi}.
	\end{align*}
Now we apply the substitution $t_k' = t_k'/t$ which gives us
	\begin{align*}
		&t^{2n(b+r-1)}\int_{\R^{nd}} \left|\frac{1}{n!} \sum_{\sigma \in \Sigma_n} \int_{[0,t]^n<} \overline{\mathcal{F}G\left(1-\frac{t_n}{t},\cdot\right)(t^{b/a}(\sum_{i=1}^n\xi_i))\cdots \mathcal{F}G\left(\frac{t_2}{t}-\frac{t_1}{t},\cdot\right)(t^{b/a}\xi_1)} \right|^2 \varphi(\vec{\xi})\ud \vec{\xi}.
		\\& = t^{2n(b+r-1)+2n}\int_{\R^{nd}} \left| \frac{1}{n!} \sum_{\sigma \in \Sigma_n} \int_{[0,1]^n<} \overline{\mathcal{F}G\left(1-t_n,\cdot\right)(t^{b/a}(\sum_{i=1}^n\xi_i))\cdots \mathcal{F}G\left(t_2-t_1,\cdot\right)(t^{b/a}\xi_1)} \right|^2 \varphi(\vec{\xi})\ud \vec{\xi}.
	\end{align*}
Next, we recall the scaling property of $\varphi$ and apply the substitution $\xi_k' = t^{b/a}\xi_k$ which yields
	\begin{align*}
		& t^{2n(b+r-1)+2n} \int_{\R^{nd}}  \left| \frac{1}{n!} \sum_{\sigma \in \Sigma_n} \int_{[0,1]^n<} \overline{\mathcal{F}G\left(1-t_n,\cdot\right)(t^{b/a}(\sum_{i=1}^n\xi_i))\cdots \mathcal{F}G\left(t_2-t_1,\cdot\right)(t^{b/a}\xi_1)} \right|^2 \varphi(\vec{\xi})\ud \vec{\xi}
		\\&= t^{2n(b+r-1)+2n}\int_{\R^{nd}} \left| \frac{1}{n!} \sum_{\sigma \in \Sigma_n} \int_{[0,1]^n<} \overline{\mathcal{F}G\left(1-t_n,\cdot\right)(\sum_{i=1}^n\xi_i)\cdots \mathcal{F}G\left(t_2-t_1,\cdot\right)(\xi_1)} \right|^2 \varphi(t^{-b/a}\vec{\xi})t^{-nd b / a}\ud \vec{\xi}
	\end{align*}

	\begin{align*}
		\\& = t^{2n(b+r-1)+2n - n b \alpha / a}\int_{\R^{nd}} \left| \frac{1}{n!} \sum_{\sigma \in \Sigma_n} \int_{[0,1]^n<} \overline{\mathcal{F}G\left(1-t_n,\cdot\right)(\sum_{i=1}^n\xi_i)\cdots \mathcal{F}G\left(t_2-t_1,\cdot\right)(\xi_1)} \right|^2 \varphi(\vec{\xi})\ud \vec{\xi}
		\\& = t^{2n(b+r-1)+2n - n b\alpha / a}\int_{\R^{nd}} |\mathcal{F}\tilde{f}_n(\cdot,0,1)|^2(\vec{\xi}) \varphi(\vec{\xi})\ud \vec{\xi}
		\\&=  t^{[2(b+r-1)+2 - b \alpha / a]n}\Norm{\tilde{f_n}(\cdot,0,1)}_{H^{\otimes n}}^2
		\\&= t^{[2(b+r) - b \alpha / a]n}\Norm{\tilde{f_n}(\cdot,0,1)}_{H^{\otimes n}}^2
	\end{align*}
$\hfill \qed$

The following result gives the existence and uniqueness of the solution. Its proof follows by a standard procedure.
	\begin{theorem}
		Fix any $T \in (0,\infty]$. Suppose that $f_n(\cdot,x,t) \in \mathcal{H}^{\otimes n}$ for any $t \in (0,T)$, $x \in R^d$ and $n \ge 1$. Then equation \ref{spdee} has a unique $L^2(\Omega)$ solution on $(0,T) \times \R^d$ if and only if the series \ref{solwexp} converges in $L^2(\Omega)$ for any $(t,x) \in (0,T) \times \R^d$, which is equivalent to the convergence of the series \ref{2mom}. In this case, the solution is given by \ref{solwexp} with the second moment given by
			\begin{equation} \label{2mom}
				\mathbb{E}(|u(t,x)|^2) = \sum_{n \ge 0} \theta^n n! \Norm{\tilde{f_n}(\cdot,x,t)}_{\mathcal{H}^{\otimes n}}^2 \quad \forall (t,x) \in (0,T) \times \R^d.
			\end{equation}
	\end{theorem}

\subsection{Some large deviation results.}\label{largedevresults} 
First we define 
	\[
		\rho := \rho_{\nu,a} = \sup_{\Norm{f}_2 =1} \int_{\R^d} \left[ \int_{\R^d} \frac{f(x+y)f(y)}{\sqrt{1+\frac{\nu}{2}|x+y|^a } \sqrt{1+\frac{\nu}{2}|y|^a} } \ud y \right]^2 \mu(\ud x)
	\]
Recall that $\mu(\ud \xi) = C_{d,\alpha}|\xi|^{-(d-\alpha)} \ud \xi$ which corresponds to the Riesz kernel $\gamma(x) = |x|^{-\alpha}$ for some $\alpha \in (0,d)$. In X. Chen \cite{chenlarge}, the case of $a \in (0,2]$ and $\nu = 2$ is discussed and there it is proven that
	\begin{equation} \label{rhowavedev}
		\lim_{n \to \infty} \frac{1}{n} \log \left[ \frac{1}{(n!)^2} \int_{(\R^d)^n} \left( \sum_{\sigma \in \Sigma_n} \prod_{k=1}^n \frac{1}{1+ |\sum_{j=k}^n \xi_{\sigma(j)}|^a}\right)^2 \mu (\ud \vec{\xi})\right] = \log[\rho_{2,a}].
	\end{equation}
and in addition, X. Chen \cite{chenlarge} also shows the equality 
	\begin{equation}\label{oldrhovarrep}
		\rho_{2,a} = \Lambda_{\alpha}^{2-(\alpha/a)}
	\end{equation}
 where
	\begin{align}
		\Lambda_{\alpha} &= \sup_{g \in \mathcal{F}_a} \left\{ \left( \int_{\R^{2d}} \frac{g^2(x_1)g^2(x_2)}{|x_1+x_2|^\alpha} \ud \vec{x} \right)^{1/2} - \mathcal{E}_a(g,g) \right\}
		\\&= \sup_{g \in \mathcal{F}_a} \left\{ \langle g^2 * g^2, |\cdot|^{-\alpha} \rangle_{L^2(\R^d)} - \mathcal{E}_a(g,g) \right\} \nonumber.
	\end{align}
and 
	\begin{align*}
		\mathcal{E}_a(g,g) := (2\pi)^{-d} \int_{\R^d} |\xi|^a |\mathcal{F}g(\xi)|^2 \ud \xi
	\end{align*}
The following result shows the relationship between $\rho_{2,a}$ and $\rho$ and well as the generalization of \ref{rhowavedev} for $\rho.$

	\begin{theorem}\label{rhovariationaltheorem}
For any $\nu >0 $ and $a \in (0,2]$, we have that
		\begin{equation}\label{rhorelation}
					\rho = \left(\frac{\nu}{2}\right)^{-\alpha/a}  \Lambda_{\alpha}^{2-(\alpha/a)}.
		\end{equation}
In addition,
		\begin{equation}\label{rhofracdev}
				\lim_{n \to \infty} \frac{1}{n} \log \left[ \frac{1}{(n!)^2} \int_{(\R^d)^n} \left( \sum_{\sigma \in \Sigma_n} \prod_{k=1}^n \frac{1}{1+ \frac{\nu}{2}|\sum_{j=k}^n \xi_{\sigma(j)}|^a}\right)^2 \mu (\ud \vec{\xi})\right] = \log[\rho].
		\end{equation}
	\end{theorem}
\noindent \textbf{Proof:} Recall that we define $\varphi(x) = C_{d,\alpha}|x|^{-(d-\alpha)}$. We begin by applying the substitution $x' = \left(\frac{\nu}{2}\right)^{1/a}$ to $\rho$:
	\begin{align*}
		\rho &= \sup_{\Norm{f}_2 =1} \int_{\R^d} \left[ \int_{\R^d} \frac{f(x+y)f(y)}{\sqrt{1+\frac{\nu}{2}|x+y|^a } \sqrt{1+\frac{\nu}{2}|y|^a} } \ud y \right]^2 \mu(\ud x)
		\\& = \sup_{\Norm{f}_2 =1} \int_{\R^d} \left[ \int_{\R^d} \frac{f\left(\left(\frac{\nu}{2}\right)^{-1/a}(x+y)\right)f\left(\left(\frac{\nu}{2}\right)^{-1/a}y\right)}{\sqrt{1+|x+y|^a } \sqrt{1+|y|^a} } \left(\frac{\nu}{2}\right)^{-d/a} \ud y \right]^2 \left(\frac{\nu}{2}\right)^{(d-\alpha)/a}\varphi(x)\left(\frac{\nu}{2}\right)^{-d/a}\ud x
		\\&= \left( \frac{\nu}{2} \right)^{-\alpha/a} \sup_{\Norm{f}_2 =1} \int_{\R^d} \left[ \int_{\R^d} \frac{f\left(\left(\frac{\nu}{2}\right)^{-1/a}(x+y)\right)f\left(\left(\frac{\nu}{2}\right)^{-1/a}y\right)}{\sqrt{1+|x+y|^a } \sqrt{1+|y|^a} } \left(\frac{\nu}{2}\right)^{-d/a} \ud y \right]^2 \varphi(x)\ud x
	\end{align*}
Now define $f^*(x) =\left(\frac{\nu}{2}\right)^{\frac{-d}{2a}} f\left( \left(\frac{\nu}{2}\right)^{-1/a} x \right)$. Then with this we see that
	\begin{align*}
		\rho &= \left( \frac{\nu}{2} \right)^{-\alpha/a} \sup_{\Norm{f}_2 =1} \int_{\R^d} \left[ \int_{\R^d} \frac{f^*\left(x+y\right)f^*\left(y\right)}{\sqrt{1+|x+y|^a } \sqrt{1+|y|^a} }  \ud y \right]^2 \varphi(x)\ud x
	\end{align*}
and note that
	\[
		\int_{\R^d} f(x)^2 \ud x = \int_{\R^d} f^*(x)^2 \ud x
	\]
and thus
	\[
		\left\{ f \in L^2(\R^d) : \Norm{f}_2 = 1 \right\} = 	\left\{ f \in L^2(\R^d) : \Norm{f^*}_2 = 1 \right\}
	\]
so we can say that
	\begin{align*}
		\rho &= \left( \frac{\nu}{2} \right)^{-\alpha/a} \sup_{\Norm{f}_2 =1} \int_{\R^d} \left[ \int_{\R^d} \frac{f^*\left(x+y\right)f^*\left(y\right)}{\sqrt{1+|x+y|^a } \sqrt{1+|y|^a} }  \ud y \right]^2 \varphi(x)\ud x
		\\&= \left( \frac{\nu}{2} \right)^{-\alpha/a}  \sup_{\Norm{f^*}_2 =1}  \int_{\R^d} \left[ \int_{\R^d} \frac{f^*\left(x+y\right)f^*\left(y\right)}{\sqrt{1+|x+y|^a } \sqrt{1+|y|^a} }  \ud y \right]^2 \varphi(x)\ud x
		\\&\left( \frac{\nu}{2} \right)^{-\alpha/a}  \rho_{2,a}
	\end{align*}
and \ref{rhorelation} is proven by applying \ref{oldrhovarrep}.

To prove \ref{rhofracdev}, we first notice that by applying the substitution $\xi_{\sigma(j)}' = \left(\frac{\nu}{2}\right)^{1/a} \xi_{\sigma(j)}$ that
	\begin{equation}
			\int_{(\R^d)^n} \left[\sum_{\sigma \in \Sigma_n} \prod_{k=1}^n \frac{1}{1+\frac{\nu}{2}\left|\sum_{j=k}^n \xi_{\sigma(j)} \right|^a}\right]^2 \mu(\ud \vec{\xi})  = \left( \frac{\nu}{2} \right)^{-n\alpha/a}  \int_{(\R^d)^n} \left[\sum_{\sigma \in \Sigma_n} \prod_{k=1}^n \frac{1}{1+\left|\sum_{j=k}^n \xi_{\sigma(j)} \right|^a}\right]^2 \mu(\ud \vec{\xi} )
	\end{equation}
and now we can directly calculate
	\begin{align*}
		\lim_{n \to \infty} \frac{1}{n} & \log \left[ \frac{1}{(n!)^2} \int_{(\R^d)^n} \left( \sum_{\sigma \in \Sigma_n} \prod_{k=1}^n \frac{1}{1+ \frac{\nu}{2}|\sum_{j=k}^n \xi_{\sigma(j)}|^a}\right)^2 \mu (\ud \vec{\xi})\right]
		\\&= \lim_{n \to \infty} \frac{1}{n} \log \left[ \frac{1}{(n!)^2}  \left( \frac{\nu}{2} \right)^{-n\alpha/a}  \int_{(\R^d)^n} \left[\sum_{\sigma \in \Sigma_n} \prod_{k=1}^n \frac{1}{1+\left|\sum_{j=k}^n \xi_{\sigma(j)} \right|^a}\right]^2 \mu(\ud \vec{\xi} )\right]
		\\&= \log\left[ \left( \frac{\nu}{2} \right)^{-\alpha/a} \right] +  \lim_{n \to \infty} \frac{1}{n} \log \left[ \frac{1}{(n!)^2}   \int_{(\R^d)^n} \left[\sum_{\sigma \in \Sigma_n} \prod_{k=1}^n \frac{1}{1+\left|\sum_{j=k}^n \xi_{\sigma(j)} \right|^a}\right]^2 \mu(\ud \vec{\xi} )\right]
	\end{align*}
and because of both Theorem 1.5 \cite{chenlarge} and \ref{rhorelation} above,
	\begin{align*}
		\lim_{n \to \infty} \frac{1}{n} \log \left[ \frac{1}{(n!)^2} \int_{(\R^d)^n} \left( \sum_{\sigma \in \Sigma_n} \prod_{k=1}^n \frac{1}{1+ \frac{\nu}{2}|\sum_{j=k}^n \xi_{\sigma(j)}|^a}\right)^2 \mu (\ud \vec{\xi})\right]  = \log\left[ \left( \frac{\nu}{2} \right)^{-\alpha/a} \right] + \log[\rho_{2,a}] = \log[\rho]
	\end{align*}
$\hfill \qed$

Next we show that $\rho$ is finite. We do this by showing $\rho_{2,a}$ is finite and then noting that this implies that $\rho$ is finite by \ref{rhovariationaltheorem}. We stated in the introduction that Chen, X. proved in \cite{chenlarge} that $\rho$ is finite for the case of the Riesz kernal. We mimic his proof but were he uses the Hardy-Littlewood-Sobelov theorem, we make use of the weak Young's inequality instead. 

	\begin{lemma}\label{rhofinite}
For any $f,g,h$ with $h\ge0$ and for $\varphi$ given as in equation \ref{specfunction}, we have the following:
\[
\left( \int_{\R^d} \left[ \int_{\R^d} \dfrac{|f(x+y)g(y)|}{\sqrt{h(x+y)}\sqrt{h(y)}} \ud y \right]^p \varphi(x) \ud x \right)^{1/p} \le C \Norm{f}_2 \Norm{g}_2 \Norm{h^{-1}}_{pd/\alpha}
\] 
where $\alpha = \sum_{=1}^k \alpha_i$ with $k \in \{1, \cdots ,d\}$, $\alpha_i \in (0,d_i)$ and $\varphi(x) = \prod_{i =1}^k |x_{(i)}|^{-(d_i-\alpha_i)}$ where $x_{(i)} = (x_{i_1},\cdots , x_{i_{d_i}} )$ is the $i^{th}$ group of $x=(x_1, \cdots, x_d)$.
	\end{lemma}
\noindent \textbf{Proof:} By observing the proof of Lemma 1.6 \cite{chenlarge}, we notice that we only need to prove equation (1.27) which says 
\begin{equation}\label{fracrhofinite}
\int_{\R^d}\int_{\R^d} F(y)G(x)\varphi(x-y) \ud y \ud x \le C \Norm{F}_{2d/(d+\alpha)} \Norm{G}_{2d/(d+ \alpha)}.
\end{equation}
where 
\[
F(x) = \frac{|f(x)|}{(h(x))^{p/2}} \quad \text{and} \quad G(x) = \frac{|g(x)|}{(h(x))^{p/2}} 
\]
If we can do this then the result follows by copying the rest of Chen's proof. To prove \ref{fracrhofinite}, we use the weak version of Young's inequality. Recall that weak Young's inequality \cite{LiebAnalysis} says
\[
\left|	\int_{\R^d}\int_{\R^d}a(x)b(x-y)c(y) \ud x \ud y \right| \le K_{p,q,r,d} \Norm{a}_p\Norm{b}_{q,w}\Norm{c}_r
\]
where $p,q,r >1$ and $1/p + 1/q + 1/r =2$ and 
\[
\Norm{b}_{q,w} = \sup_{A} |A|^{-1/q'} \int_A |f(x)| \ud x
\]
where $A$ is an arbitrary Borel set of finite measure. So to prove \ref{fracrhofinite} we wish to be able to apply this with
\[
p=r=2d/(d+\alpha) \quad \text{and} \quad q= d/\left(d-\sum_{j=1}^d\alpha_j \right) = d/(d-\alpha)
\]
and, in addition, with 
\[  
a=F, \quad c=G \quad \text{and} \quad b=\varphi
\]
Again, by following the steps of the proof from the Xia Chen paper, we note that it suffices to show that $\Norm{\varphi}_{q,w}$ is finite with $q = d/(d-\alpha)$. Note that $1/q' = 1 - 1/q = \alpha/d.$

First we define $A_R := A_1 \times \cdots \times A_k$ where $A_i \subset \R^{d_i}$ and $A_i = [-R,R]$ if $d_i = 1$ and $A_i = B_R(0)$ if $d_i > 1$ where $B_R(0)$ is the ball centered at the origin with radius $R$. With this we have that 
\begin{equation}\label{int} \int_{A_i} |x_{(i)}|^{-(d_i-\alpha_i)} \ud x_{(i)} =
\begin{cases} 
2 \dfrac{R^{\alpha_i}}{\alpha_i} & , d_i =1
\\ |\mathbb{S}^{d_i - 1}|\dfrac{R^{\alpha_i}}{\alpha_i} & , d_i >1
\end{cases}
\end{equation}
where we used polar coordinated to calculate the integral when $d_i > 1$ and $|\mathbb{S}^{d_i-1}|$ is the surface area of the sphere in $\R^{d_i}$. Also, 
\begin{equation}\label{leb} |A_i|  =
\begin{cases}
2 R& , d_i =1
\\ |\mathbb{S}^{d_i - 1}|\dfrac{R^{d_i}}{d_i} & , d_i >1.
\end{cases}
\end{equation}
Recall that $1/q' = \alpha/d$. Then by combining \ref{int} and \ref{leb} we see that 
\begin{align*}
|A_R|^{-1/q'}	\int_{A_R} \varphi(x) \ud x &= \left(\prod_{i =1}^k |A_i|\right)^{-1/q'} \prod_{i=1}^k \int_{A_i} |x_{(i)}|^{-(d_i-\alpha_i)} \ud x_{(i)}
\\&= R^{-\alpha} \left( \prod_{i =1}^k\begin{cases}
2 & , d_i =1
\\ |\mathbb{S}^{d_i - 1}| / d_i & , d_i >1.
\end{cases} \right) \times R^\alpha 	\left(\prod_{i =1}^k \begin{cases} 
 2/\alpha_i & , d_i =1
\\ |\mathbb{S}^{d_i - 1}| / \alpha_i & , d_i >1
\end{cases} \right)
\\&= K
\end{align*}
where $K$ is a constant independent of $R$. Hence we that 
\[
\sup_{R>0} 	|A_R|^{-1/q'}	\int_{A_R} \varphi(x) \ud x = K.
\]
and by the symmetry of $\varphi$ we can in fact say that 
\[
\sup_{B} 	|B|^{-1/q'}	\int_{B} \varphi(x) \ud x = K
\]
where $B$ denotes the Borel sets of finite measure. Hence, $\varphi \in L_{q,w}(\R^d)$ with $q = d/(d-\alpha)$ which completes the proof.

$\hfill \qed$

We finish this section by making a slight generalization of $\Lambda_{\alpha}$. For this, we define 
	\begin{equation}
		\mathcal{M}_a(\Theta f , \theta) = \sup_{g \in \mathcal{F}_a} \left\{ \langle g^2 * g^2, f \rangle^{1/2} - \frac{\theta}{2}\mathcal{E}_a(g,g) \right\}.
	\end{equation}
One notices that by making this generalization, we can apply Theorem \ref{mainresult} to the SPDE \ref{spdee} driven by an arbitrary Gaussian noise with correlation function $f$ as long as one can show that $\rho$ is finite. By slightly modifying the proof of Lemma 2.3 \cite{BalanChenChen}, one can show that 
	\begin{equation}
			\mathcal{M}_a(\Theta f , \theta) = \Theta^{\frac{a}{2a-\alpha}}\theta^{-\frac{\alpha}{2a-\alpha}} \mathcal{M}_a( f , 1).
	\end{equation}
Finally we note that 
	\begin{equation}
		\Lambda_{\alpha} = \mathcal{M}_a(|\cdot|^{-\alpha},2) = 2^{-\frac{\alpha}{2a-\alpha}}\mathcal{M}_a(|\cdot|^{-\alpha},1) 
	\end{equation}
and by defining $\mathcal{M}_a := \mathcal{M}_a(|\cdot|^{-\alpha},1) $ and by applying Theorem \ref{rhovariationaltheorem} we see that 
	\begin{equation}\label{rhoMrepresentation}
		\rho = \left(\frac{\nu}{2}\right)^{-\alpha/a}    \left[2^{-\frac{\alpha}{2a-\alpha}}\mathcal{M}_a\right]^{2-(\alpha/a)} = \nu^{-\alpha/a} \mathcal{M}_a^{\frac{2a-\alpha}{a}}
	\end{equation}
	
\section{Solvability}\label{solvability}
The goal for this section is to prove the following theorem:
	\begin{theorem}\label{thmsolve}
		The equation \ref{spdee} has a unique solution $\{u(t,x) : t>0, x\in \R^d\}$ and
			\begin{equation}
				\sup_{(t,x) \in [0,T] \times \R^d} \mathbb{E}(|u(t,x)|^p) < \infty \quad \text{for all } p \ge 2, T>0
			\end{equation}
	\end{theorem}

We need to build up some theory before we prove the above theorem. We first introduce some notation that will be used throughout the remainder of his paper. We start by defining $L_n(\vec{y})$ to be the Laplace transform of $\tilde{f_n}(\vec{y},0,\cdot)$ evaluated at $1$:
	\begin{align*}
		L_n(\vec{y}) &= n! \int_0^\infty e^{-t}\tilde{f_n}(\vec{y};0,t) \ud t
		\\&= \sum_{\sigma \in \Sigma_n} \int_0^\infty e^{-t} \int_{[0,t]^n<} \prod_{k=1}^n G(s_k - s_{k-1}, y_{\sigma(k)}-y_{\sigma(k-1)}) \ud \vec{s} \ud t
	\end{align*}
with the convention that. $s_0 = y_{\sigma(0)} = 0$. Next we define,
	\begin{align*}
		H_n(t, \vec{x}) &= n! \int_{(\R^d)^n} \prod_{k=1}^n K(x_k - y_k) \tilde{f_n}(\vec{y};0,t) \ud \vec{y}
		\\&= \sum_{\sigma \in \Sigma_n} \int_{[0,t]^n <} \int_{(\R^d)^n} \prod_{k=1}^n K(x_k-y_k) 	\prod_{k=1}^n G(s_k - s_{k-1}, y_{\sigma(k)} - y_{\sigma(k-1)}) \ud \vec{y} \ud \vec{s}
	\end{align*}
where in the last line we used relation \ref{fperm0} for expressing  $\tilde{f_n}(\cdot;0,t)$. Next, similarly to $L_n$, we define $L(x)$ and can calculate $\mathcal{F}L(\xi)$ to be
	\begin{equation} \label{LandFL}
		L(x) = \int_0^\infty e^{-t}G(t,x) \ud t \quad \text{and} \quad \mathcal{F}L(\xi) = \int_0^\infty e^{-t}\mathcal{F}G(t,\cdot)(\xi) \ud t = \dfrac{1}{1 + \frac{\nu}{2}|\xi|^a}
	\end{equation}
and through a change of variables, one sees that the Laplace transform of convolution in the time variable become a product:
	\begin{equation}
			L_n(\vec{y}) = \sum_{\sigma \in \Sigma_n} L(y_{\sigma(1)}) L(y_{\sigma(2)}-y_{\sigma(1)}) \cdots L(y_{\sigma(n)}-y_{\sigma(n-1)})
	\end{equation}
and using this and \ref{LandFL}, we can calculate $\mathcal{F}L_n(\vec{\xi})$:
	\begin{equation}\label{FLfn}
		\mathcal{F}L_n(\vec{\xi}) = \sum_{\sigma \in \Sigma_n} \prod_{k=1}^{n} \dfrac{1}{1+ \frac{\nu}{2}\left|\sum_{j=k}^n \xi_j \right|^a}
	\end{equation}

	\begin{lemma}\label{Hneq1}
Under the assumptions of Theorem \ref{mainresult}, we have that
					\begin{equation}
						\int_{(\R^d)^n} \left[ \int_0^\infty e^{-t} H_n(t,\vec{x}) \ud t \right]^2 \ud \vec{x} = \int_{(\R^d)^n} | \mathcal{F}L_n(\vec{\xi}) |^2 \mu(\ud \vec{\xi}).
					\end{equation}
	\end{lemma}
\noindent\textbf{Proof:} 
	\begin{align*}
		 \int_0^\infty e^{-t} H_n(t,\vec{x}) \ud t &= \sum_{\sigma \in \Sigma_n} \int_0^\infty e^{-t} \int_{[0,t]^n <} \int_{(\R^d)^n} \prod_{k=1}^n K(x_k-y_k) 	\prod_{k=1}^n G(s_k - s_{k-1}, y_{\sigma(k)} - y_{\sigma(k-1)}) \ud \vec{y} \ud \vec{s} \ud t
		 \\&= \int_{\R^{dn}} \prod_{k=1}^n K(x_k-y_k) L(\vec{y}) \ud \vec{y}.
	\end{align*}
Therefore,
	\begin{align*}
		\int_{\R^{dn}}\left[\int_0^\infty e^{-t} H_n(t,\vec{x}) \ud t \right]^n \ud \vec{x} &= 	\int_{\R^{dn}} \left[  \int_{\R^{dn}} \prod_{k=1}^n K(x_k-y_k) L(\vec{y}) \ud \vec{y} \right]^2 \ud \vec{x} 
		\\&= \int_{\R^{dn}} |\mathcal{F}L_n(\vec{\xi})|^2 \prod_{k=1}^n\mu(\ud \xi_k) 
	\end{align*}
where the last line follows from Plancherel's theorem and the fact that $K*K =\gamma$. \textcolor{red}{consider adding the plancherel details}	

$\hfill \qed$
	\begin{lemma}\label{inteq1}
Under the assumptions of Theorem \ref{mainresult}, we have that
					\begin{equation}
						\Norm{\tilde{f_n}(\cdot,0;t)}_{\mathcal{H}^{\otimes n}}^2 = \frac{1}{(n!)^2}\int_{(\R^d)^n} H_n^2(t,\vec{x}) \ud \vec{x}.
					\end{equation}
	\end{lemma}
\noindent\textbf{Proof:} Writing $\gamma(y_k-y_k') = \int_{\R^d} K(y_k-x_k)K(y_k'-x_k) \ud x_k$, we see that 
	\begin{align*}
		\Norm{\tilde{f_n}(\cdot,0;t)}_{\mathcal{H}^{\otimes n}}^2 &= \int_{(\R^d)^n} \int_{(\R^d)^n} \prod_{k=1}^n \gamma(y_k-y_k') \tilde{f}_n(\vec{y},0;t)\tilde{f}_n(\vec{y'},0;t) \ud \vec{y} \ud \vec{y'}
		\\&= \int_{(\R^d)^n} \left[ \int_{(\R^d)^n} \prod_{k=1}^n K(y_k-x_k) \tilde{f}_n(\vec{y},0;t) \ud \vec{y} \right]^2 \ud (\vec{x})
		\\&=  \int_{(\R^d)^n} \left[ \frac{1}{n!} H_n(t,\vec{x}) \right]^2 \ud \vec{x}
	\end{align*}
\textcolor{red}{maybe add the placherel theorem steps in the above equalities.}

$\hfill \qed$

If $f$ is a non-decreasing function on $[0,\infty)$, then the following reverse Cauchy-Schwarz inequality holds, which plays a key role in our developments.

	\begin{lemma}\label{intineq1}
		If $f:[0,\infty) \to [0,\infty)$ is a non-decreasing function, then
			\begin{equation}
				2\int_0^\infty e^{-2t}f^2(t) \ud t \le \left( \int_0^\infty e^{-t} f(t) \ud t \right)^2
			\end{equation}
	\end{lemma}
\noindent\textbf{Proof:} Let $\tau$ and $\tau'$ be two independent random variables with exponential distribution of mean 1. Then $\tau \wedge \tau'$ has an exponential distribution with density $g(t)=2e^{-2t}$ and
	\[
		\left( \int_0^\infty e^{-t}f(t) \ud t \right)^2 = \mathbb{E}[ f(\tau) ] \mathbb{E}[ f(\tau') ] =  \mathbb{E}[ f(\tau) f(\tau')] = 2 \int_0^\infty e^{-2t} f^2(t) \ud t
	\]
$\hfill \qed$

Denote
	\begin{equation}
		T_n = T_n(\nu,a) := \int_{(\R^d)^n} \left[\sum_{\sigma \in \Sigma_n} \prod_{k=1}^n \frac{1}{1+\frac{\nu}{2}\left|\sum_{j=k}^n \xi_{\sigma(j)} \right|^a}\right]^2 \mu(\ud \vec{\xi}).
	\end{equation}
By applying the substitution $\xi_k' = \left(\frac{\nu}{2}\right)^{1/a}\xi_k$ and the scaling property for $\varphi$, we may rewrite $T_n$ as
	\begin{equation}
		T_n = \left( \frac{\nu}{2} \right)^{-n\alpha/a} T_n(2,a)
	\end{equation}

	\begin{lemma}
		Under the assumption of Theorem \ref{momentasym}, we have that
			\begin{equation}\label{muassumption}
				C_\mu = C_\mu(\nu,a) := \int_{(\R^d)} \left(\frac{1}{1+\frac{\nu}{2}|\xi|^a}\right)^2 \mu(\ud \xi) < \infty \quad \text{and} \quad T_n \le (n!)^2(C_\mu)^2
			\end{equation}
	\end{lemma}
\noindent \textbf{Proof:}
As in the proof of Lemma 3.6 \cite{BalanChenChen}, we apply a spherical substitution and get that
	\[
		\int_{(\R^d)} \left(\frac{1}{1+|\xi|^a}\right)^2 \mu(\ud \xi) = \left(\int_{S^{d-1}}\varphi(\gamma) \ud \sigma(\gamma)\right) \left( \int_0^\infty \frac{\rho^{\alpha-1}}{(1+\rho^{a})^2} \ud \rho \right).
	\]
The first integral is shown to be finite in \cite{BalanChenChen}. The second integral is finite if
	\begin{align}\label{alpha2}
		2a - \alpha +1 > 1
	\end{align}
which is already satisfied under the assumptions of Theorem \ref{mainresult}. Hence the first relation in \ref{muassumption} is proven.

The second relation is proven the same was as in Lemma 3.6 \cite{BalanChenChen}. Here is the proof. 
	\begin{align}
		T_n &\le n! \int_{(\R^d)^n} \sum_{\sigma \in \Sigma_n}\left[ \prod_{k=1}^n \frac{1}{1+\frac{\nu}{2}\left|\sum_{j=k}^n \xi_{\sigma(j)} \right|^a}\right]^2 \mu(\ud \xi_1)\cdots \mu(\ud \xi_n) \nonumber
		\\ &= (n!)^2 \int_{(\R^d)^n} \left[ \prod_{k=1}^n \frac{1}{1+\frac{\nu}{2}\left|\sum_{j=k}^n \xi_j \right|^a}\right]^2 \mu(\ud \xi_1)\cdots \mu(\ud \xi_n) \nonumber
		\\& \le (n!)^2 \left[ \sup_{\eta \in \R^d} \int_{(\R^d)^n}  \prod_{k=1}^n \left( \frac{1}{1+\frac{\nu}{2}\left| \xi + \eta \right|^a}\right)^2 \mu(\ud \xi) \right]^n \nonumber
		\\& = (n!)^2 \left[ \int_{(\R^d)^n}  \prod_{k=1}^n \left(\frac{1}{1+\frac{\nu}{2}\left| \xi \right|^a}\right)^2 \mu(\ud \xi) \right]^n \nonumber
		\\&= (n!)^2(C_\mu)^n \label{lem62step}
	\end{align}
where the last line follows by Lemma 4.1 of Balan and Song \cite{BalanSong} or by Lemma 3.1 in \cite{Chen2007}.

$\hfill \qed$

The following lemma shows that $\tilde{f}_n(\cdot,0;t) \in \mathcal{H}^{\otimes n}$ for any $t>0$.
	\begin{lemma}\label{LNfnineq}
Under the assumptions of Theorem \ref{mainresult} and for any $t>0$.
			\[
				\int_0^\infty e^{-t} \Norm{\tilde{f_n}(\cdot,0;t)}_{ \mathcal{H}^{ \otimes n } }^n \ud t \le \dfrac{2^{n\left[2(b+r)-\frac{b\alpha}{a}\right]}}{(n!)^2} T_n \le 2^{n\left[2(b+r)-\frac{b\alpha}{a}\right]} (C_\mu')^n
			\]
	\end{lemma}

\noindent \textbf{Proof:} We first derive the following scaling property for $\mathcal{F}\tilde{f}_n$ where $\tilde{f}_n$ is the kernel of the $n^{th}$ Weiner integral in the chaos expansion of the solution \ref{solwexp}:
	\begin{equation}\label{Ffnscale}
		\mathcal{F}\tilde{f}_n(\cdot,0,ct)(\xi_1,\cdots,\xi_n) = c^{n(b+r)} \mathcal{F}\tilde{f}_n(\cdot,0,t)(c^{b/a}\xi_1, \cdots, c^{b/a}\xi_n).
	\end{equation}
Let $c>0$.
	\begin{align*}
		\mathcal{F}\tilde{f}_n(\cdot,0,ct)(\xi_1,\cdots,\xi_n) = \frac{1}{n!} \sum_{\sigma \in \Sigma_n} \mathcal{F}f_n(\cdot,0,ct)(\xi_{\sigma(1)},\cdots, \xi_{\sigma(n)})
	\end{align*}
so it suffices to just look at the $\mathcal{F}f(\cdot,0,ct)(\xi_{\sigma(1)},\cdots, \xi_{\sigma(n)})$ terms individually.
	\begin{align*}
		\mathcal{F}f_n(\cdot,0,ct)(\xi_1,\cdots, \xi_n) &= \int_{[0,ct]^n_<} \prod_{k=1}^n \overline{ \mathcal{F}G(t_{k+1} - t_k, \cdot)\left( \sum_{j=1}^k \xi_j \right) }  \ud \vec{t}
		\\&= \int_{[0,t]^n_<} \prod_{k=1}^n \overline{ \mathcal{F}G\left(c(t_{k+1} - t_k), \cdot \right)\left( \sum_{j=1}^k \xi_j \right) } c^n \ud \vec{t}
	\end{align*}
where we applied the substitution $t_i' = t_i/c$ by applying the scaling property of $\mathcal{F}G$ we wee that
	\begin{align*}
		\int_{[0,t]^n_<} \prod_{k=1}^n  \overline{ \mathcal{F}G\left(c(t_{k+1} - t_k), \cdot \right)\left( \sum_{j=1}^k \xi_j \right) } c^n \ud \vec{t} &= \int_{[0,t]^n_<} \prod_{k=1}^n \overline{ c^{b+r-1} \mathcal{F}G\left(t_{k+1} - t_k, \cdot \right)\left( c^{b/a} \sum_{j=1}^k \xi_j \right) } c^n \ud \vec{t}
		\\&= c^{n(b+r)} \mathcal{F}f_n(\cdot,0,t)(c^{b/a}\xi_1, \cdots, c^{b/a}\xi_n)
	\end{align*}
Now by using this scaling property we see that
	\begin{align*}
		\int_0^\infty e^{-t} \Norm{ \tilde{f}_n(\cdot,0,t) }_{\mathcal{H}^{\otimes n}}^2 \ud t &= \int_0^\infty e^{-t} \int_{\R^{nd}} \left| \mathcal{F}\tilde{f_n}(\cdot,0;t)(\xi_1,\cdots,\xi_n) \right|^2 \mu(\ud \vec{\xi})  \ud t
		\\&= \int_0^\infty e^{-2t} \int_{\R^{nd}} \left| \mathcal{F}\tilde{f_n}(\cdot,0;2t)(\xi_1,\cdots,\xi_n) \right|^2 \mu(\ud \vec{\xi})  2 \ud t
		\\& = 2^{2n(b+r)}\int_0^\infty 2e^{-2t} \int_{\R^{nd}} \left| \mathcal{F}\tilde{f_n}(\cdot,0,t)(2^{b/a}\xi_1,\cdots,2^{b/a}\xi_n) \right|^2 \mu(\ud \vec{\xi})   \ud t
		\\& = 2^{2n(b+r)}\int_0^\infty 2e^{-2t} \int_{\R^{nd}} \left| \mathcal{F}\tilde{f_n}(\cdot,0;t)(\xi_1,\cdots,\xi_n) \right|^2 2^{\frac{-nbd}{a}} 2^{\frac{nb(d-\alpha)}{a}} \mu(\ud \vec{\xi})   \ud t
		\\& = 2^{n(2(b+r) - b\alpha/a)}  \int_0^\infty 2e^{-2t} \int_{\R^{nd}} \left| \mathcal{F}\tilde{f_n}(\cdot,0;t)(\xi_1,\cdots,\xi_n) \right|^2 \mu(\ud \vec{\xi})   \ud t
		\\&=  \frac{2^{n(2(b+r) - b\alpha/a)}}{(n!)^2}  \int_0^\infty 2e^{-2t} \int_{\R^{nd}} H_n(t,\vec{x})^2 \ud \vec{x}   \ud t
	\end{align*}
and now we apply Lemmas \ref{inteq1} and \ref{intineq1} as well as Equation \ref{FLfn} to get
	\begin{align*}
		\frac{2^{n(2(b+r) - b\alpha/a)}}{(n!)^2}  \int_0^\infty 2e^{-2t} \int_{\R^{nd}} H_n(t,\vec{x})^2 \ud \vec{x}   \ud t  & \le \frac{2^{n(2(b+r) - b\alpha/a)}}{(n!)^2}  \int_{\R^{nd}} \left(\int_0^\infty e^{-t}  H_n(t,\vec{x}) \ud t \right)^2 \ud \vec{x}
		\\&=   \frac{2^{n(2(b+r) - b\alpha/a)}}{(n!)^2}  \int_{\R^{nd}} | \mathcal{F}L_n(\vec{\xi}) |^2\ud \vec{\xi}
		\\&= \frac{2^{n(2(b+r) - b\alpha/a)}}{(n!)^2}  T_n
	\end{align*}
 and this proves the first inequality. An application of Lemma 3.6 gives is the second inequality.

 $\hfill \qed$

\noindent\textbf{Proof of Theorem \ref{thmsolve}:} First recall \ref{Ffnscale}, which says
	\[
		\mathcal{F}\tilde{f_n}(\cdot,0,ct)(\xi_1,..., \xi_n) = c^{n(b+r)}\mathcal{F}\tilde{f_n}(\cdot,0,t)(c^{b/a}\xi_1,...,c^{b/a}\xi_n).
	\]

In order to prove the theorem we need to show that $\sum_{n \ge 0} \theta^n n! \Norm{ \tilde{f_n}(\cdot,0,t) }_{ \mathcal{H}^{\otimes n} }^2 < \infty.$ From the scaling property \ref{normscale} we see that
	\begin{align*}
		\int_0^\infty e^{-t} \Norm{ \tilde{f_n}(\cdot,0,t) }_{ \mathcal{H}^{\otimes n} }^2 \ud t & = \Norm{ \tilde{f_n}(\cdot,0,1) }_{ \mathcal{H}^{\otimes n} }^2 \int_0^\infty e^{-t} t^{ \left[ 2(b+r)-\frac{b \alpha}{a} \right] n } \ud t
		\\&= \Norm{ \tilde{f_n}(\cdot,0,1) }_{ \mathcal{H}^{\otimes n} }^2 \Gamma\left( \left[ 2(b+r)- \frac{b \alpha}{a} \right] n +1 \right)
	\end{align*}
as long as we have
	\begin{equation} \label{alpha3}
		2(b+r)-\frac{b\alpha}{a} >0
	\end{equation}
By using Lemma \ref{LNfnineq} we can deduce that
	\begin{align*}
		\Norm{ \tilde{f_n}(\cdot,0,1) }_{ \mathcal{H}^{\otimes n} }^2 &= \dfrac{ \int_0^\infty e^{-t} \Norm{ \tilde{f_n}(\cdot,0,t) }_{ \mathcal{H}^{\otimes n} }^2 \ud t }{ \Gamma\left( \left[ 2(b+r)- \frac{b \alpha}{a} \right] n +1 \right) }
		\\ &\le \dfrac{ 2^{n \left[ 2(b+r)- \frac{b \alpha}{a} \right]} }{ \Gamma\left( \left[ 2(b+r)- \frac{b \alpha}{a} \right] n +1 \right) } T_n
		\\& \le \dfrac{ \left(2^{ \left[ 2(b+r)- \frac{b \alpha}{a} \right]}C_\mu' \right)^n }{ \Gamma\left( \left[ 2(b+r)- \frac{b \alpha}{a} \right] n +1 \right) }
	\end{align*}

By the same reasoning as in the proof in the wave paper (details in my notebook) we can say that there exists a constant depending on $\alpha$, $k_\alpha$, such that
	\[
		\Norm{ \tilde{f_n}(\cdot,0,1) }_{ \mathcal{H}^{\otimes n} }^2 \le k_{\alpha} \dfrac{ \left(2^{ \left[ 2(b+r)- \frac{b \alpha}{a} \right]}C_\mu' \right)^n }{ (n!)^{ 2(b+r) - b\alpha/a } }
	\]

Therefore,
	\begin{align*}
		\sum_{n \ge 0} \theta^n n! \Norm{ \tilde{f_n}(\cdot,0,t) }_{ \mathcal{H}^{\otimes n} }^2 & \le \sum_{n \ge 0} \theta^n n!  k_{\alpha} \dfrac{ \left(2^{ \left[ 2(b+r)- \frac{b \alpha}{a} \right]}C_\mu' \right)^n }{ (n!)^{ 2(b+r) - b\alpha/a } }
		\\& =  \sum_{n \ge 0} \theta^n   k_{\alpha} \dfrac{ \left(2^{ \left[ 2(b+r)- \frac{b \alpha}{a} \right]}C_\mu' \right)^n }{ (n!)^{ 2(b+r) - 1 - b\alpha/a } }
	\end{align*}
and the last sum is finite when
	\begin{equation}\label{alpha4}
		2(b+r) - 1 - b\alpha/a > 0
	\end{equation}

$\hfill \qed$

\section{The case $p=2$.} In this section, we will prove the following result.
	\begin{theorem}\label{p=2}
Under the assumptions of Theorem \ref{mainresult}, we have that 
			\begin{equation}
			\lim_{t \to \infty}  \frac{1}{t^\beta} \log \mathbb{E}(|u(t,x)|^2) \nonumber = \left( \frac{2a}{2a(b+r)- b\alpha} \right)^\beta (\theta\rho)^{\frac{a}{2a(b+r)-b\alpha-a}}\left(2(b+r)-\frac{b\alpha}{a}-1\right).
			\end{equation}
	\end{theorem}

In order to prove this theorem, we first need to prove the following lemma.

	\begin{lemma}
		\begin{align}
			\lim_{n \to \infty}  \frac{1}{n} \log \left( \Norm{ \tilde{f_n}(\cdot,0,1) }_{ \mathcal{H}^{\otimes n} }^2 \Gamma\left( \left[ 2(b+r)- \frac{b \alpha}{a} \right] n +1 \right) \right) &= \log\left(2^{2(b+r)-\frac{b\alpha}{a}}\right) + \log\left[\rho\right] \label{2ineqa}
		\end{align}

		\begin{align}
			\lim_{n \to \infty}  \frac{1}{n} \log \left( (n!)^{2(b+r)- \frac{b \alpha}{a} }\Norm{ \tilde{f_n}(\cdot,0,1) }_{ \mathcal{H}^{\otimes n} }^2  \right) &=  \log\left[\left(\frac{2}{2(b+r)-\frac{b\alpha}{a}}\right)^{2(b+r)-\frac{b\alpha}{a}}\right] + \log\left[\rho\right] \label{2ineqb}
		\end{align}
	\end{lemma}

\noindent \textbf{Proof:} Recall the definition of $T_n$ from above as well as \ref{rhofracdev} which says that
	\[
			\lim_{n \to \infty} \frac{1}{n} \log \left[ \frac{1}{(n!)^2} \int_{(\R^d)^n} \left( \sum_{\sigma \in \Sigma_n} \prod_{k=1}^n \frac{1}{1+ \frac{\nu}{2}|\sum_{j=k}^n \xi_{\sigma(j)}|^a}\right)^2 \mu (\ud \vec{\xi})\right] = \lim_{n \to \infty} \frac{1}{n} \log \left[ \frac{T_n}{(n!)^2} \right]  = \log[\rho].
	\]
We already showed in the proof of Theorem \ref{thmsolve} above that
	\[
		\int_0^\infty e^{-t} \Norm{ \tilde{f_n}(\cdot,0,t) }_{ \mathcal{H}^{\otimes n} }^2 \ud t = \Norm{ \tilde{f_n}(\cdot,0,1) }_{ \mathcal{H}^{\otimes n} }^2 \Gamma\left( \left[ 2(b+r)- \frac{b \alpha}{a} \right] n +1 \right)
	\]
and then Lemma \ref{LNfnineq} above shows that
	\[
		\Norm{ \tilde{f_n}(\cdot,0,1) }_{ \mathcal{H}^{\otimes n} }^2 \Gamma\left( \left[ 2(b+r)- \frac{b \alpha}{a} \right] n +1 \right) \le \dfrac{2^{n\left[2(b+r)-\frac{b\alpha}{a}\right]}}{(n!)^2} T_n
	\]
Combining this and \ref{rhofracdev} we see that
	\begin{align*}
		\limsup_{n \to \infty} \frac{1}{n}& \log\left[ \Norm{ \tilde{f_n}(\cdot,0,1) }_{ \mathcal{H}^{\otimes n} }^2 \Gamma\left( \left[ 2(b+r)- \frac{b \alpha}{a} \right] n +1 \right)  \right]
		\\	&\le \log\left( 2^{\left[2(b+r)-\frac{b\alpha}{a}\right]} \right) + \lim_{n \to \infty} \frac{1}{n} \log\left( \frac{T_n}{(n!)^2} \right)
		\\& = \log\left(  2^{\left[2(b+r)-\frac{b\alpha}{a}\right]} \right)  + \log(\rho)
	\end{align*}
which proves the upper bound for \ref{2ineqa}.

Now we prove the lower bound. Define for $t>0$ and $t'>0$
	\[
		J_n(t,t') = \int_{\R^{nd}} H_n(t,x)H_n(t',x) \ud x
	\]
and recall Lemma \ref{inteq1} which says
	\[
		\Norm{\tilde{f_n}(\cdot,0;t)}_{\mathcal{H}^{\otimes n}}^2 = \frac{1}{(n!)^2} \int_{(\R^{d})^n} H_n(t,\vec{x})^2\ud \vec{x} = \frac{1}{(n!)^2} J_n(t,t).
	\]
Thus $J_n(t,t)$ satisfies the scaling property
	\[
		J_n(t,t) =  t^{[2(b + r ) -b \alpha / a]n} J_n(1,1).
	\]
Let $\tau$ and $\tilde{\tau}$ be independent exponential random variables with mean 1. By lemma \ref{Hneq1},
	\begin{align*}
		\mathbb{E}[J_n(\tau, \tau')] &= \int_0^\infty \int_0^\infty e^{-t}e^{-\tilde{t}}J_n(t,\tilde{t}) \ud t \ud \tilde{t} = \int_{\R^{dn}}\left[ \int_0^\infty e^{-t}H_n(t,x) \ud t \right]^2 \ud x
		\\&= \int_{\R^{nd}} |\mathcal{F}L_n(\xi_1,...,\xi_n)|^2 \mu(\ud \xi_1) \cdots \mu(\ud \xi_n) = T_n
	\end{align*}
In addition,
	\[
		J_n(t,t') \le J_n(t,t)^{1/2}J_n(t',t')^{1/2} =  t^{[2(b + r ) -b \alpha / a]n/2} {t'}^{[2(b + r ) -b \alpha / a]n/2} J_n(1,1).
	\]
Combining all of this gives us
	\begin{align*}
		T_n & = \mathbb{E}[J_n(\tau,\tilde{\tau})] \le \mathbb{E}\left[  {\tau}^{[2(b + r ) -b \alpha / a]n/2} \right] \mathbb{E}\left[  {\tilde{\tau}}^{[2(b + r ) -b \alpha / a]n/2} \right]J_n(1,1)
		\\& = \Gamma\left( \frac{[2(b + r ) -b \alpha / a]n}{2} \right)^2 J_n(1,1)
		\\& = \Gamma\left( \frac{[2(b + r ) -b \alpha / a]n}{2} \right)^2 (n!)^2	\Norm{ \tilde{f_n}(\cdot,0;1) }_{ \mathcal{H}^{\otimes n} }^2
	\end{align*}
where the last line follows from Lemma \ref{inteq1}. Thus so far we have shown
	\begin{align}
		\frac{T_n}{(n!)^2} = \frac{\left( \frac{\nu}{2} \right)^{-n\alpha/a} T_n'}{(n!)^2}\le	\Gamma\left( \frac{[2(b + r ) -b \alpha / a]n}{2} \right)^2 	\Norm{ \tilde{f_n}(\cdot,0;1) }_{ \mathcal{H}^{\otimes n} }^2 \label{2lowbd}
	\end{align}
Now, another application of Stirling's formula shows that
	\[
		\Gamma\left( \frac{[2(b + r ) -b \alpha / a]n}{2} \right)^2  \sim \Gamma\left( [2(b + r ) -b \alpha / a]n +1 \right)2^{- [2(b + r ) -b \alpha / a]n}C_n
	\]
	where $C_n = 2^{-1}[2(b + r ) -b \alpha / a]^{1/2}(2 \pi n )^{1/2}$. We now infer from \ref{2lowbd} that
	\[
		c_{\alpha} \frac{2^{[2(b + r ) -b \alpha / a]n}}{C_n} \frac{1}{(n!)^2} T_n \le \Gamma\left( [2(b + r ) -b \alpha / a]n +1 \right) 	\Norm{ \tilde{f_n}(\cdot,0,1) }_{ \mathcal{H}^{\otimes n} }^2
	\]
where $c_{\alpha}$ is a constant depending on $\alpha$. Since $\frac{1}{n}\log(C_n) \to 0$ as $n \to \infty$, we obtain the lower bound
	\begin{align*}
		\liminf_{n\to \infty} \frac{1}{n} &\log \left( \Gamma\left( [2(b + r ) -b \alpha / a]n +1 	\right) 	\Norm{ \tilde{f_n}(\cdot,0,1) }_{ \mathcal{H}^{\otimes n} }^2  \right)
		\\ & \ge \log\left(  2^{[2(b + r ) -b \alpha / a]} 	\right) + \log(\rho)
	\end{align*}
Thus \ref{2ineqa} is proven. Lastly, \ref{2ineqb} follows from \ref{2ineqa} and the fact that
	\[
		\lim_{n \to \infty} \frac{1}{n} \log\left[ \frac{\Gamma(an+1)}{(n!)^a} \right] = a\log(a).
	\]
	
$\hfill \qed$

With the above lemma, we can now prove Theorem \ref{p=2}.

\noindent \textbf{Proof of Theorem \ref{p=2}:} Assuming that there is an $L^2(\Omega)$ solution $u(t,x)$ and then applying the scaling property \ref{normscale} gives us
	\begin{align*}
		\mathbb{E}(|u(t,x)|^2)  &= \sum_{n\ge 0} 	\theta^n(n!)\Norm{\tilde{f}_n(\cdot,0,t)}_{\mathcal{H}^{\otimes n}}^2 = \sum_{n \ge 0}\theta^n(n!)t^{(2(b+r)-b\alpha/a)n}\Norm{\tilde{f}_n(\cdot,0,1)}_{\mathcal{H}^{\otimes n}}^2
		\\&= \sum_{n \ge 0}x_nR_nt^{(2(b+r)-b\alpha/a)n}
	\end{align*}
where
	\[
		R_n = (n!)^{2(b+r)-b\alpha/a}\Norm{\tilde{f}_n(\cdot,0,1)}_{\mathcal{H}^{\otimes n}}^2 \quad \text{and} \quad x_n=\frac{\theta^n}{(n!)^{2(b+r)-(b\alpha/a)-1}}
	\]
\ref{2ineqb} above says that
	\[
		\frac{1}{n}\log(R_n) \to  \log\left[ \left(\frac{2}{2(b+r)-\frac{b\alpha}{a}}\right)^{2(b+r)-\frac{b\alpha}{a}}  \rho \right]  \quad \text{as } n\to \infty.
	\]
Now define $R$ to be
	\[
		R=\left(\frac{2}{2(b+r)-\frac{b\alpha}{a}}\right)^{2(b+r)-\frac{b\alpha}{a}}  \rho.
	\]
We want to find a $\beta$ and $A$ so that
	\[
		\lim_{t \to \infty}\frac{1}{t^\beta} \log \sum_{n \ge 0} x_nR^n\left(t^{(2(b+r)-b\alpha/a)}\right)^n = A
	\]
By lemma A3 notice that
	\[
		\lim_{t\to \infty} \left[\frac{1}{(\theta R)t^{2(b+r)-b\alpha/a}}\right]^{\frac{1}{2(b+r)-(b\alpha/a)-1}} \log \sum_{n \ge 0} \frac{\left[ (\theta R) t^{2(b+r)-b\alpha/a}\right]^n}{(n!)^{2(b+r)-(b\alpha/a)-1}} = 2(b+r) - (b\alpha/a) -1
	\]
and by an easy algebraic manipulation
	\begin{align*}
		\lim_{t\to \infty} \left[\frac{1}{t^{2(b+r)-b\alpha/a}}\right]^{\frac{1}{2(b+r)-(b\alpha/a)-1}} & \log \sum_{n \ge 0} \frac{\left[ (\theta R) t^{2(b+r)-b\alpha/a}\right]^n}{(n!)^{2(b+r)-(b\alpha/a)-1}}
		\\&= [2(b+r) - (b\alpha/a) -1](\theta R)^{\frac{1}{2(b+r)-(b\alpha/a)-1}}
	\end{align*}
hence we see that
	\[
		\beta = \frac{2(b+r)-b\alpha/a}{2(b+r)-(b\alpha/a)-1} \quad \text{and} \quad A=[2(b+r) - (b\alpha/a) -1](\theta R)^{\frac{1}{2(b+r)-(b\alpha/a)-1}}
	\]
or by removing the fractions in the denominator
	\[
		\beta = \frac{2a(b+r)-b\alpha}{2a(b+r)-(b\alpha)-a} \quad \text{and} \quad A=[2(b+r) - (b\alpha/a) -1](\theta R)^{\frac{a}{2a(b+r)-b\alpha-a}}
	\]
Finally, lemma A2 says that
	\begin{align}
		\lim_{t\to \infty} & \frac{1}{t^\beta} \log \mathbb{E}(|u(t,x)|^2) \nonumber
		\\&= \theta^{\frac{a}{2a(b+r)-b\alpha-a}}\left( \frac{2a}{2a(b+r)- b\alpha} \right)^\beta \rho^{\frac{a}{2a(b+r)-b\alpha-a}}\left(2(b+r)-\frac{b\alpha}{a}-1\right)
	\end{align}

\section{Upper Bound.}\label{upperbound} In this section we use the result obtained in Theorem \ref{p=2} and the hypercontractivity property given by Theorem B.1 in \cite{BalanChenChen} to prove the following theorem.

	\begin{theorem}
Under the assumptions of Theorem \ref{mainresult}
		\begin{equation}
			\limsup_{n \to \infty} \; t_p^{-\beta} \log \Norm{u(t,x)}_p  \le \left(\frac{1}{2}\right)  \left( \frac{2a}{2a(b+r)- b\alpha} \right)^\beta (\theta\rho)^{\frac{a}{2a(b+r)-b\alpha-a}}\left(2(b+r)-\frac{b\alpha}{a}-1\right).
		\end{equation}
	\end{theorem}
\noindent \textbf{Proof:} We denote $u_{\lambda}$ for $\lambda >0$ to be the the solution of the fractional equation with $\theta$ replaced with $\lambda$ and $u=u_\theta$. Using theorem B1 from the wave paper we can say that for any $p \ge 2$ that
	\[
		\Norm{u(t,0)}_p \le \Norm{u_{(p-1)\theta}(t,0)}_2
	\]
and by the scaling property
	\begin{align*}
		\Norm{u_{(p-1)\theta}(t,0)}_2^2 &= \sum_{n \ge 0}(n!)^n\theta^n(p-1)^n \Norm{\tilde{f}_n(\cdot,0,t)}_{\mathcal{H}^{\otimes n}}^2
		\\&= \sum_{n \ge 0}(n!)^n\theta^n \Norm{\tilde{f}_n\left(\cdot,0,t(p-1)^{\frac{1}{2(b+r)-b\alpha/a}}\right)}_{\mathcal{H}^{\otimes n}}^2
		\\&= \Norm{u\left(t(p-1)^{\frac{1}{2(b+r)-b\alpha/a}}, 0\right)}_2^2
	\end{align*}
and from this we can see that
	\[
		\Norm{u(t,0)}_p \le \Norm{u\left(t(p-1)^{\frac{1}{2(b+r)-b\alpha/a}}, 0\right)}_2 =: \Norm{u(t_p,0)}_2
	\]
and an application of Theorem \ref{p=2} proves the result.

$\hfill \qed$

\section{Lower Bound.}\label{lowerbound} The goal of this section is to prove the following theorem.

	\begin{theorem}\label{lbound}
Under the assumptions of Theorem \ref{mainresult}
			\begin{equation}
				\liminf_{n \to \infty} t_p^{-\beta} \log \Norm{u(t,x)}_p \ge \left(\frac{1}{2}\right) \nonumber \left( \frac{2a}{2a(b+r)- b\alpha} \right)^\beta 		(\theta \rho)^{\frac{a}{2a(b+r)-b\alpha-a}}\left(2(b+r)-\frac{b\alpha}{a}-1\right).
			\end{equation}
	\end{theorem}

We will need to build up some theory before being able to prove the above theorem. We start by defining the function $W_n(t,\phi)$ on $(0,\infty) \times L_{\mathbb{C}}^2(\mu)$ by
	\begin{equation}
		W_n(t,\phi) := \int_{[0,t]^n<}
		\int_{(\R^d)^n} \prod_{k=1}^n \phi(\xi_k) \prod_{k=1}\mathcal{F}G(s_k-s_{k-1},\cdot)(\xi_k + \cdots + \xi_n) \mu(\ud \xi_1)\cdots \mu(\ud \xi_n) \ud s
	\end{equation}
with $s_0=0$. We now give conditions for when $W_n$ is well defined.

	\begin{lemma}\label{LWn}
If the measure $\mu$ satisfies the relation in \ref{muassumption}, then $W_n(t,\phi)$ is well defined and for any $d \ge 1$, $t>0$ and $\phi \in L^2_{\mathbb{C}}(\mu)$. Moreover, 
		\begin{equation}\label{lem62}
			\int_0^\infty e^{-t}W_n(t,\phi) \ud t = \int_{(\R^d)^n} \prod_{k=1}^n \phi(\xi_k) \prod_{k=1}^n \frac{1}{1+\frac{\nu}{2}|\xi_k + \cdots + \xi_n|^a} \mu(\ud \xi) \cdots \mu(\ud \xi_n)
		\end{equation}
	\end{lemma}

\noindent \textbf{Proof:} By Fubini's theorem, the left hand side of \ref{lem62} is equal to
	\begin{align*}
		 \int_{(\R^d)^n} \prod_{k=1}^n \phi(\xi_k) \int_0^\infty e^{-t} \int_{[0,t]^n_<} \prod_{k=1}^n \mathcal{F}G(s_k-s_{k-1}, \cdot)(\xi_k + \cdots + \xi_n) \ud \vec{s} \ud t \mu(\ud \vec {\xi}),
	\end{align*}
which is equal to, by the change of variables $t-s_n = u$ and $s_k - s_{k-1} = u_k$ for $k=1,...,n$, 
	\begin{align*}
		 \int_{(\R^d)^n} \prod_{k=1}^n \phi(\xi_k) \left( \int_0^\infty e^{-u} \ud u \right) \prod_{k=1}^n \left( \int_0^\infty e^{-u_k} \mathcal{F}G(u_k,\cdot)(\vec{\xi}) \; \ud u_k \right) \prod_{k=1}^n \mu(\ud \xi_k).
	\end{align*}
Now recall that $\int_0^\infty e^{-t} \mathcal{F}G(t,\cdot)(\xi) \; \ud t = 1/(1+\frac{\nu}{2}|\xi|^a)$. Then the above quantity is equal to 
	\[
		 \int_{(\R^d)^n} \prod_{k=1}^n \phi(\xi_k) \prod_{k=1}^n \frac{1}{1+|\sum_{j=k}^n \xi_j|^2} \prod_{k=1}^n \mu(\ud \xi_k).
	\]
The result it follows from applying the Cauchy-Schwarz inequality. In addition, by using the same arguments that lead to \ref{lem62step}, we see the above quantity is finite. 

$\hfill \qed$ 

	\begin{proposition}\label{Wnprop1} 
For $f\in \mathcal{H}$, $t>0$, and $p,q>1$ with $p^{-1} + q^{-1}=1$, it holds that
		\[
			\Norm{u(t,0)}_p \ge \exp\left\{ -\frac{1}{2}(q-1)\Norm{f}_{\mathcal{H}}^2 \right\} \left| \sum_{n \ge 0} \theta^{n/2} W_n(t,\mathcal{F}f) \right|
		\]
and as a consequence, the series $\left| \sum_{n \ge 0} \theta^{n/2} W_n(t,\mathcal{F}f) \right|$ converges provided that $\Norm{u(t,0)}_p < \infty$, which is the case under Theorem \ref{thmsolve}.
	\end{proposition}

\noindent \textbf{Proof:} Suppose $f \in \mathcal{H}$ and let 
	\[
		Z_f = \exp\left( W(f) - \frac{1}{2}\Norm{f}_\mathcal{H}^2 \right) = \sum_{n \ge 0} \frac{1}{n!} I_n(f^{\otimes n}).
	\]
If $f$ is a distribution, then $f^{\otimes n}$ is the distribution in $\mathcal{S}'(\R^{nd})$ whose Fourier transform is given by: $\mathcal{F}f^{\otimes n}(\vec{\xi}) = \prod_{k=1}^n \mathcal{F}f(\xi_k)$. If $f$ is a function, then $f^{\otimes n}$ is also a function and is defined by $f^{\otimes n}(\vec{x}) = \prod_{k=1}^n f(x_k)$. Now using Holder's inequality and  the fact that $W(f) \sim N(0,\Norm{f}_\mathcal{H}^2)$, we obtain:
	\begin{equation}\label{Zfholder1}
			|\mathbb{E}\left[ u(t,0)Z_f \right]| \le \Norm{u(t,0)}_p \Norm{Z_f}_q = \Norm{u(t,0)}_p  \exp\left\{ \frac{1}{2}(q-1)\Norm{f}_{\mathcal{H}}^2 \right\}.
	\end{equation}
Now by the orthogonality of the Wiener Chaos spaces,
	\begin{equation}\label{Zfholder2}
		|\mathbb{E}\left[ u(t,0)Z_f \right]| \le \sum_{n \ge 0} \theta^{n/2} \frac{1}{n!} \mathbb{E}[I_n(f_n(\cdot,0;t)) I_n(f^{\otimes n} )] =  \sum_{n \ge 0} \theta^{n/2} \langle \tilde{f}_n(\cdot,0;t), f^{\otimes n} \rangle_{\mathcal{H}^{\otimes n} } .
	\end{equation}
By using the definition of $\tilde{f}_n(\cdot,0;t)$, we see that 
	\begin{align*}
		 \langle \tilde{f}_n(\cdot,0;t), f^{\otimes n} \rangle_{\mathcal{H}^{\otimes n} } &= \int_{(\R^d)^n} \prod_{k=1}^n \mathcal{F}f(\xi_k) \overline{ \mathcal{F}\tilde{f}_n (\cdot, 0 ; t)(\vec{\xi} ) } \prod_{k=1}^n \mu(\ud \xi_k)
		 \\&= \frac{1}{n!} \sum_{\sigma \in \Sigma_n} \int_{[0,t]^n_<} \ud \vec{t} \int_{(\R^{d})^n} \prod_{k=1} \mu(\ud \xi_k) \prod_{k=1} \mathcal{F}f(\xi_k) \prod_{k=1} \mathcal{F}G(t_{k+1}-t_k,\cdot)\left( \sum_{k=1}^n \xi_{\sigma(k)} \right)
		 \\&= \int_{[0,t]^n_<} \ud \vec{t} \int_{(\R^{d})^n} \prod_{k=1}^n \mu(\ud \xi_k) \prod_{k=1}^n \mathcal{F}f(\xi_k) \prod_{k=1}^n \mathcal{F}G(t_{k+1}-t_k,\cdot)\left( \sum_{k=1}^n \xi_k \right)
		 \\&= W_n(t,\mathcal{F}f)
	\end{align*}
where $t_{n+1} = t$ and the last step is due to the change of variables $s_1 = t - t_n, \cdots, s_n = t - t_1$ and $\xi_1' = \xi_n, \cdots, \xi_n' = \xi_1$. The conclusion follows from the above centered equations and \ref{Zfholder1} and \ref{Zfholder2}.

$\hfill \qed$

Our next goal is to prove the following proposition, which is where the motivation of the definitions of $t_p$ and $\beta$ come from.

	\begin{proposition}\label{Wnprop2}
 For $p,q > 1$ and $p^{-1}+q^{-1}=1$ we have that
		\[
			\Norm{u(t,0)}_p \ge \exp\left\{ -\frac{1}{2} t_p^\beta \Norm{f}_{\mathcal{H}}^2 \right\}\left| \sum_{n \ge 0} \theta^{n/2} W_n(t_p^\beta, \mathcal{F}f) \right|
		\]
	\end{proposition}

\noindent \textbf{Proof:} The above inequality motivates us to find a function $f_{\tau}(x)$ and some $\tau$ such that
	\begin{enumerate}
		\item \[(q-1)\Norm{f_{\tau}}_{\mathcal{H}}^2 = (p-1)^{-1}\Norm{f_{\tau}}_{\mathcal{H}}^2 = t_p^\beta \Norm{f}_{\mathcal{H}}^2\]
		\item  \[W_n(t, \mathcal{F}f_\tau) = W_n(t_p^\beta, \mathcal{F}f)\]
	\end{enumerate}
since then we could just apply Proposition \ref{Wnprop1} to the function $f_\tau$. Note above that we have not yet defined $\beta$ or $t_p$. We guess $f_\tau$ takes the form $f_\tau(x) = \tau^Vf(\tau^Wx)$. We now mention, through direct calculation one can show that $\Norm{f_\tau}_\mathcal{H}^2$ and $W_n(t,\mathcal{F}f_\tau)$ can be calculated as follows:
	\begin{equation}
		\Norm{f_\tau}_{\mathcal{H}}^2 = \tau^{2(V-dW)+W\alpha} \Norm{f}_\mathcal{H}^2
	\end{equation}
and
	\begin{equation}
		W_n (t,\mathcal{F}f_\tau) =	\tau^{n\left[ V-W\left( (d-\alpha)+ \frac{a}{b}(b+r) \right) \right]}W_n(t\tau^{\frac{a}{b}W}, \mathcal{F}f).
	\end{equation}
This means we must have
	\[
		\begin{cases}
			V-W\left( (d-\alpha)+ \frac{a}{b}(b+r) \right) &= 0
			\\ (p-1)^{-1}\tau^{2(V-dW)+W\alpha} &= t\tau^{\frac{a}{b}W}
		\end{cases}.
	\]
We need to now choose an appropriate $\tau$ and we choose $\tau = (p-1)t$ and so the above system of equations becomes
	\[
		\begin{cases}
			V-W\left( (d-\alpha)+ \frac{a}{b}(b+r) \right) &= 0
			\\ (p-1)^{2(V-dW)+W\alpha-1}t^{2(V-dW)+W\alpha} &= (p-1)^{\frac{a}{b}W}t^{\frac{a}{b}W+1}
		\end{cases}.
	\]
Equating the exponents of either the $(p-1)$ term or the $t$ term both lead to the following equation
	\begin{equation}
		2(V-dW)+W\alpha-1 - \frac{a}{b}W = 0
	\end{equation}
and so now we have two linear equations in $W$ and $V$ that we can solve, namely
	\[
		\begin{cases}
			V-W\left( (d-\alpha)+ \frac{a}{b}(b+r) \right) &= 0
			\\2(V-dW)+W\alpha-1 - \frac{a}{b}W &= 0
		\end{cases}.
	\]
Now if we define
	\begin{equation}
		\beta := \dfrac{2(b+r)-\dfrac{b\alpha}{a}}{ 2(b+r)-\dfrac{b\alpha}{a} -1 }
	\end{equation}
then solving the system gives us
	\[
		\begin{cases}
			W = \dfrac{b}{a}\left( \beta -1 \right)
			\\V = \left(\frac{a}{b}(b+r)-\alpha +d\right)\frac{b}{a}(\beta -1)
		\end{cases}.
	\]
Lastly if we define
	\begin{equation}
		t_p^\beta := (q-1)(p-1)^\beta t^\beta = (p-1)^{\beta -1} t^\beta
	\end{equation}
then Proposition \ref{Wnprop2} is proven by applying Proposition \ref{Wnprop1} to the function
	\[
		f_\tau(x) = \tau^{\left(\frac{a}{b}(b+r)-\alpha+d\right)\frac{b}{a}(\beta -1)} f\left(\tau^{\frac{b}{a}(\beta -1)}x\right)
	\]
with $\tau=(p-1)t$.

$\hfill \qed$

To remove the absolute value sign in Proposition \ref{Wnprop2}, we want to identify all the $f \in \mathcal H$ for which $W_n(t,\mathcal{F}f)$ is nonnegative. In fact, if we consider the space
	\[
		\mathcal{H}_+ = \{ f\in \mathcal{H} : \text{$f$ is nonegative and nonnegative definite} \}
	\]
then by Plancherel's theorem
	\[
			W_n(t,\mathcal{F}f) = \int_{[0,t]^n<} \int_{\R^{nd}} \prod_{k=1}^n(f*\gamma)(x_k) \prod_{k=1} G(s_k - s_{k-1}, x_k-x_{k-1}) \ud \vec{x} \ud \vec{s} =: U_n(t,f)
	\]
with $s_0=0$ and $x_0=0$. Under Assumption \ref{paramassumptions}, the Green's function for the fractional equation is nonnegative and hence under these conditions we get that $W_n(t,\mathcal{F}(f)) \ge 0$ for all the functions $f \in \mathcal{H}_+$. Hence we have that 
	\[
		0 \le W_n(t, \mathcal{F}f) = U_n(t,f) < \infty \quad \text{for all } f \in \mathcal{H}_+ \text{ and for any of the four options in Aummption \ref{paramassumptions}.}
	\]
where the finiteness comes from Proposition \ref{Wnprop2}.

Now we define
	\[
		W_n(t) := \sup_{f\in\mathcal{H}, \Norm{f}_{\mathcal{H}}=1} W_n(t,\mathcal{F}(f)) \quad \text{and} \quad U_n(t) :=  \sup_{f\in\mathcal{H_+}, \Norm{f}_{\mathcal{H}}=1} U_n(t,f)
	\]
and with this we see that
	\[
		W_n(t) \ge U_n(t) \ge 0
	\]

	\begin{theorem}
If $\tau$ is an exponential random variable with mean $1$, then
		\[
			\liminf_{n\to \infty} \frac{1}{n} \log \mathbb{E} (U_n(\tau) ) \ge \log(\rho^{1/2})
		\]
	\end{theorem}
\noindent \textbf{Proof:} We start by letting $\tau$ be an exponential random variable with mean $1$. With this,
	\[
		\mathbb{E}(U_u(\tau)) = \int_0^\infty e^{-t} U_n(t) \ud t \ge \sup_{f\in \mathcal{H_+}n \Norm{f}_\mathcal{H} =1} \int_0^\infty e^{-t} U_n(t,f) \ud t
	\]
For an $f \in \mathcal{H}_+$ and $\Norm{f}_{\mathcal{H}}= 1$ we have that $\mathcal{F}f$ is non-negative and non-negative definite which further implies that $\mathcal{F}f$ is even. In addition, Lemma \ref{LWn} gives us that
	\[
		\int_0^{\infty} e^{-t} U_n(t,f) \ud t = \int_0^\infty e^{-t} W_n(t,\mathcal{F}f) \ud t = \int_{\R^{dn}} \prod_{k=1}^n \mathcal{F}f(\xi_k) \prod_{k=1}^n \frac{1}{1+ \frac{\nu}{2}|\xi_k+\cdots + \xi_n|^a} \mu(\ud \xi_1) \cdots \mu(\ud \xi_n)
	\]
Modifying the ideas in the proof of Theorem 1.2 \cite{chenlarge} gives us the following identity:
	\begin{align*}
		\lim_{n\to \infty}\frac{1}{n} \log \int_{\R^{dn}} \prod_{k=1}^n(\mathcal{F}f)(\xi_k) \prod_{k=1}^n\frac{1}{1+\frac{\nu}{2}|\xi_k+\cdots+\xi_n|^a} \mu(\ud \xi_1) \cdots \mu(\ud \xi_n) = \log[ \rho(\mathcal{F}f)]
	\end{align*}
where $\rho(\cdot)$ is defined on the space of non-negative and non-negative definite functions and is given by
	\[
		\rho(g) := \sup_{\Norm{h}_{L_2}=1} \int_{\R^d}g(\xi)\left[ \int_{\R^d} \frac{h(\xi + \eta)h(\eta)}{\sqrt{1+\frac{\nu}{2}|\xi+\eta|^a}\sqrt{1+\frac{\nu}{2}|\eta|^a}}  \ud \eta\right] \mu(\ud \xi).
	\]
Note that since $g$ is non-negative and since $\mu$ has a non-negative density $\varphi$ then we may assume $h$ is also non-negative for the remainder of this proof. With this being said, we see that for any $f \in \mathcal{H}_+$ and $\Norm{f}_H =1$
	\[
		\liminf_{n \to \infty} \frac{1}{n}\log \mathbb{E}(U_n(t)) \ge \log \rho(\mathcal{F}f)
	\]
We now need to calculate
	\[
		\sup_{f \in \mathcal{H}_+, \Norm{f}_{\mathcal{H}}=1} \rho(\mathcal{F}f)
	\]
Consider a non-negative function $h\in L^2(\R^d)$. The function $h(\cdot)/\sqrt{1+\frac{\nu}{2}|\cdot|^a} \in L^2(\R^d)$ so that $g_h(x) := (2\pi)^{d/2} \mathcal{F}^{-1}\left( \frac{h(\cdot)}{\sqrt{1+\frac{\nu}{2}|\cdot|^a}} \right)(x)$ is well defined. \textcolor{red}{I still need to verify the details that} under these conditions, $g \in W_{1,a}(\R^d)$ with $\Norm{g}_{W_{1,a}}=1$. Here is a sketch of the details. \textcolor{red}{we no longer need this part since we do not need g to be continuous, only need to to be non-negative and non-negative definite.}

	\begin{align*}
		\Norm{g}_{W_{1,a}}^2 &= \Norm{g}^2_{L^2} + (2\pi)^{-d}\sum_{k=1}^d \int_{\R^d} \xi_k^a|\mathcal{F}g(\xi)|^2 \ud \xi
		\\&= (2\pi)^{-d}\int_{\R^d}(1+|\xi|^a)|\mathcal{F}g(\xi)|^2 \ud \xi
		\\&= \int_{\R^d} h(\xi)^2 \ud \xi =1
	\end{align*}
where above we used the fact that $\Norm{g}_{L^2}^2 = (2\pi)^{-d}\Norm{\mathcal{F}g}_{L^2}^2$. 

\textcolor{red}{We need from this point on.}

Now, for notation let $\tilde{h}(x) = h(-x)$. Note that
	\begin{align*}
		\int_{\R^d} \frac{h(\xi + \eta)h(\eta)}{\sqrt{1+\frac{\nu}{2}|\xi+\eta|^a}\sqrt{1+\frac{\nu}{2}|\eta|^a}}  \ud \eta&= (2\pi)^{-d}\int_{\R^d} \mathcal{F}g_h(\xi-\eta)\mathcal{F}g_h(-\eta) \ud \eta
		\\& =(2\pi)^{-d} (\mathcal{F}g_h*\widetilde{\mathcal{F}g_h})(\xi)
	\end{align*}
Since $h(\cdot)$ is real valued then $\widetilde{\mathcal{F}g_h} = \overline{\mathcal{F}\bar{g_h}}= \mathcal{F}\bar{g_h} $. Using this and the fact that  $\mathcal{F}(fg) = (2\pi)^{-d}\mathcal{F}(f)*\mathcal{F}(g)$, we see that $(2\pi)^{-d}(\mathcal{F}g_h*\widetilde{\mathcal{F}g_h})(\xi) =\mathcal{F}\left[|g_h|^2\right](\xi) $ and with this we see that
	\begin{align*}
		\rho(\mathcal{F}f) &= \sup_{\Norm{h}_{L^2}=1}(2\pi)^{-d} \int_{\R^d}\mathcal{F}f(\xi)(\mathcal{F}g_h*\widetilde{\mathcal{F}g_h})(\xi) \mu(\ud \xi)
		\\&= \sup_{\Norm{h}_{L^2}=1} \int_{\R^d}\mathcal{F}f(\xi)\mathcal{F}[|g_h|^2](\xi) \mu(\ud \xi)
		\\&= \sup_{\Norm{h}_{L^2}=1}\langle f, |g_h|^2 \rangle _{\mathcal{H}}.
	\end{align*}
With this we see that
	\[
		\sup_{ f\in \mathcal{H}, \Norm{f}_{\mathcal{H}=1} } \rho(\mathcal{F}f) = 	\sup_{ f\in \mathcal{H}, \Norm{f}_{\mathcal{H}=1} }\sup_{\Norm{h}_{L^2}=1} \langle f, |g_h|^2 \rangle _{\mathcal{H}} = \sup_{\Norm{h}_{L^2}=1} \langle |g_h|^2, |g_h|^2 \rangle^{1/2} _{\mathcal{H}}
	\]
where the optimal $f$ is chosen to be $|g_h|^2/\Norm{|g_h|^2}_{\mathcal{H}}$. Now we want to show that 
	\begin{equation}\label{GH+}
		|g_h|^2 \in \mathcal{H}_+
	\end{equation}
which would imply that the supremum can be restricted to $\mathcal{H_+}$. If this were the case then we would have that
	\[
		\sup_{ f \in \mathcal{H}_+, \Norm{f}_{\mathcal{H}}=1 } \rho(\mathcal{F}f) = \sup_{\Norm{h}_{L^2}=1} \langle |g_h|^2, |g_h|^2 \rangle^{1/2} _{\mathcal{H}}
	\]
and notice that
	\begin{align*}
		\langle |g_h|^2, |g_h|^2 \rangle^{1/2} _{\mathcal{H}} &= \left( \int_{\R^d} \left|\mathcal{F}[|g_h|^2](\xi)\right|^2 \mu(\ud \xi) \right)^{1/2}
		\\&= (2\pi)^{-d} \left( \int_{\R^d}\left| (\mathcal{F}g_h*\widetilde{\mathcal{F}}g_h)(\xi)\right|^2 \mu(\ud \xi) \right)^{1/2}
		\\&=  \left(\int_{\R^d}\left[ \int_{\R^d} \frac{h(\xi + \eta)h(\eta)}{\sqrt{1+\frac{\nu}{2}|\xi+\eta|^a}\sqrt{1+\frac{\nu}{2}|\eta|^a}}  \ud \eta \right]^2 \mu(\ud \xi) \right)^{1/2}
		\\& = \rho^{1/2}
	\end{align*}
which finally leads us to the equality
	\[
		\liminf_{n\to \infty} \frac{1}{n} \log \mathbb{E} (U_n(\tau) ) \ge \log(\rho^{1/2}).
	\]
Hence, we finish the proof by proving \ref{GH+}. First note that from the above we see that 
	\[
		\int_{\R^d} | \mathcal{F}|g_h|^2(\xi) |^2 \mu(\ud \xi) \le \rho < \infty
	\]
and so $|g_h|^2 \in \mathcal{H}$. Lastly, since $\mathcal{F}|g_h|^2(\xi)$ is non-negative for $\xi \in \R^d$ and for non-negative $h\in L^2(\R^d)$ then $\nu(\ud \xi) = \mathcal{F}|g_h|^2(\xi) \ud \xi$ is a non-negative tempered measure. The Bochner-Schwarz theorem then implies that $|g_h|^2$ is non-negative and non-negative definite. This shows that $|g_h|^2 \in \mathcal{H}_+$. 

$\hfill \qed$

We now give scaling properties for $W_n$ and $U_n$.

	\begin{lemma}
We have that
		\[
			W_n(t) = t^{n\frac{b}{a}\left(a- \frac{\alpha}{2}+\frac{a}{b}r \right)} W_n(1)
		\]
and moreover
		\[
			U_n(t) = t^{n\frac{b}{a}\left(a- \frac{\alpha}{2}+\frac{a}{b}r \right)} U_n(1)
		\]
	\end{lemma}
\noindent \textbf{Proof:} Let $\phi = \mathcal{F}f$ and $c>0$
	\begin{align*}
		W_n & (t,\phi(c^{-1}\cdot) = \int_{[0,t]^n_<}\int_{\R^{nd}} \prod_{k=1}^n \phi(c^{-1}\xi_k) \prod_{k=1}^n G(s_k - s_{k-1},\cdot)(\xi_k + \cdots \xi_n) \prod_{i=1}^n \varphi(\xi_i)\ud \xi_i\ud s_i
		\\&= c^{nd} \int_{[0,t]^n_<}\int_{\R^{nd}} \prod_{k=1}^n \phi(\xi_k) \prod_{k=1}^n G(s_k - s_{k-1},\cdot)(c(\xi_k + \cdots \xi_n)) \prod_{i=1}^n \varphi(c\xi_i)\ud \xi_i\ud s_i
		\\&= c^{nd}c^{-n(d-\alpha)}c^{-n\frac{a}{b}(b+r-1)} \int_{[0,t]^n_<}\int_{\R^{nd}} \prod_{k=1}^n \phi(\xi_k) \prod_{k=1}^n G\left(c^{\frac{b}{a}}(s_k - s_{k-1}),\cdot\right)(\xi_k + \cdots \xi_n) \prod_{i=1}^n \varphi(\xi_i)\ud \xi_i\ud s_i
		\\&= c^{nd}c^{-n(d-\alpha)}c^{-n\frac{a}{b}(b+r-1)} c^{-n\frac{a}{b}} \int_{\left[0,c^{\frac{a}{b}}t\right]^n_<}\int_{\R^{nd}} \prod_{k=1}^n \phi(\xi_k) \prod_{k=1}^n G\left(s_k - s_{k-1},\cdot\right)(\xi_k + \cdots \xi_n) \prod_{i=1}^n \varphi(\xi_i)\ud \xi_i\ud s_i
		\\&= c^{-n\left( a- \alpha +\frac{a}{b}r \right)} W_n\left(c^{\frac{a}{b}}t,\phi \right)
	\end{align*}
Now we apply this with $c^{\frac{a}{b}} = \frac{1}{t}$ to get
	\[
		W_n\left(t,\phi(t^{\frac{b}{a}}\cdot) \right)= t^{n\frac{b}{a}\left( a- \alpha +\frac{a}{b}r \right)} W_n\left(1,\phi \right)
	\]
Now note that
	\[
		\Norm{\phi(t^{\frac{b}{a}} \cdot )}_{L^2(\mu)}^2 = t^{-\frac{\alpha b }{a}}\Norm{\phi}_{L^2(\mu)}^2
	\]
and from this we deduce that
	\[
		\Norm{t^{\frac{\alpha b }{2 a}}\phi(t^{\frac{b}{a}} \cdot )}_{L^2(\mu)}^2 =  \Norm{\phi}_{L^2(\mu)}^2 = \Norm{f}_{\mathcal{H}}^2
	\]
and in addition
	\[
		W_n\left(t,t^{\frac{b \alpha}{2 a}}\phi(t^{\frac{a}{b}}\cdot) \right)= t^{n\frac{b}{a}\left( a- \frac{\alpha}{2} +\frac{a}{b}r \right)} W_n\left(1,\phi \right)
	\]
and the result follows after taking the sup over all $f \in \mathcal{H}$ with $\Norm{f}_\mathcal{H} =1$.

$\hfill \qed$

	\begin{lemma}
		\[
			\liminf_{n\to \infty} \frac{1}{n}\log \left[ (n!)^{ \frac{b}{a}(a-\frac{\alpha}{2} + \frac{a}{b}r)  } U_n(1) \right] \ge \log\left[ \frac{\rho^{1/2}}{ \left(\frac{b}{a}(a-\frac{\alpha}{2} + \frac{a}{b}r)\right)^{\frac{b}{a}(a-\frac{\alpha}{2} + \frac{a}{b}r)} } \right]
		\]
	\end{lemma}
\noindent \textbf{Proof:} Let $\tau$ be an exponential random variable with mean $1$. By lemma 6.7,
	\[
		\mathbb{E}(U_n(\tau)) = \mathbb{E}\left( \tau^{n\frac{b}{a}(a-\frac{\alpha}{2}+\frac{a}{b}r)} \right) U_n(1) = \Gamma\left( n\frac{b}{a}\left(a-\frac{\alpha}{2} +\frac{a}{b}r \right) +1 \right) U_n(1)
	\]
and by Theorem 6.6 we obtain
	\begin{align}\label{lb1}
		\liminf_{n\to \infty} \frac{1}{n} \log \left[ \Gamma\left( n\frac{b}{a}\left(a-\frac{\alpha}{2} +\frac{a}{b}r \right) +1 \right) U_n(1) \right] \ge \log(\rho^{1/2}).
	\end{align}
Recall that
	\[
		\lim_{n \to \infty} \frac{1}{n}\log\left( \frac{\Gamma(an+1)}{(n!)^a} \right) = a \log(a)
	\]
and with this we have that
	\[
		\liminf_{n\to \infty} \frac{1}{n} \log \left[ \frac{ \Gamma\left( n\frac{b}{a}\left(a-\frac{\alpha}{2} +\frac{a}{b}r \right) +1 \right) }{ (n!)^{\frac{b}{a}\left(a-\frac{\alpha}{2} +\frac{a}{b}r \right)} } \right] = \frac{b}{a}\left(a-\frac{\alpha}{2} +\frac{a}{b}r \right) \log\left( \frac{b}{a}\left(a-\frac{\alpha}{2} +\frac{a}{b}r \right) \right)
	\]
and the result is proven after multiplying and dividing $(n!)^{\frac{b}{a}\left(a-\frac{\alpha}{2} +\frac{a}{b}r \right)} $ on the inside of the log in \ref{lb1}

$\hfill \qed$

	\begin{lemma}\label{logUnlb}
For any $k,\theta >0$, there exists a constant $c_1 = c_1(\alpha,\mathcal{M}, k, \theta) >0$ such that, by setting $n_t = [c_1 t]$, it holds that
		\begin{equation}\label{69ineq}
			\liminf_{t \to \infty} \frac{1}{t} \log( k^{n_t} \theta^{n_t/2}U_{n_t}(t)) \ge \left( k \sqrt{\theta} \right)^{\frac{2a}{2ab-\alpha b + 2ar}} \left( \rho^{1/2} \right)^{\frac{2a}{2ab-\alpha b + 2ar}}
		\end{equation}
	\end{lemma}
\noindent \textbf{Proof:} Let $\epsilon > 0$. Lemma 6.8 guarantees the existence of an $N_\epsilon >0$ so for all $n>N_\epsilon$
	\[
		(n!)^{\frac{b}{a}\left( a - \frac{\alpha}{2} + \frac{a}{b}r \right)} U_n(1) \ge \exp\left( n(\log(R)-\epsilon) \right) = R^n e^{-n\epsilon }
	\]
where
	\[
		R = \frac{ \rho^{1/2} } { \left( \frac{b}{a}\left( a - \frac{\alpha}{2} + \frac{a}{b}r \right) \right)^{\frac{b}{a}\left( a - \frac{\alpha}{2} + \frac{a}{b}r \right) }}.
	\]
Now fix a $c>0$ and let $n_t := [ct]$. Notice that $n_t \ge N_\epsilon$ for any $t>t_\epsilon := (N_\epsilon +1)/c$.

Let $t > t_\epsilon$.
	\begin{align*}
		k^{n_t}\theta^{n_t/2}U_{n_t}(t) = 	k^{n_t} \theta^{n_t/2} t^{\frac{b}{a} \left( a-\frac{\alpha}{2} + \frac{a}{b}r \right)n_t } U_{n_t}(1) \ge \frac{ k^{n_t} \theta^{n_t/2} t^{\frac{b}{a} \left( a-\frac{\alpha}{2} + \frac{a}{b}r \right)n_t } }{ (n!)^{\frac{b}{a} \left(a-\frac{\alpha}{2} +\frac{a}{b}r \right)} } R^{n_t}e^{-n_t \epsilon}
	\end{align*}
Notice that $[ct]/t \to c$ as $t \to \infty$ which means that $n_t/t \to c$ as $t \to \infty$. With this we can say
	\[
		\liminf_{t \to \infty} \frac{1}{t} \log \left( k^{n_t} \theta^{n_t/2} U_{n_t}(t) \right) = c \liminf_{t\to \infty} \frac{1}{n_t} \log \left( k^{n_t} \theta^{n_t/2} U_{n_t}(t) \right)  =: I(n_t)
	\]
Now,
	\begin{align*}
		I(n_t) & \ge  c \liminf_{t\to \infty} \frac{1}{n_t} \log \left( (kR\sqrt{\theta})^{n_t} \frac{ t^{\frac{b}{a} \left( a-\frac{\alpha}{2} + \frac{a}{b}r \right)n_t } }{ (n_t!)^{\frac{b}{a}\left( a -\frac{\alpha}{2} + \frac{a}{b}r \right) }  }\right) - c\epsilon
		\\&= c \log(k \sqrt{\theta}R) + c \liminf_{t \to \infty} \frac{1}{n_t} \log \left[  \left(\frac{t}{n_t}\right)^{\frac{b}{a}\left(a-\frac{\alpha}{2} + \frac{a}{b}r\right)n_t} \frac{ n_t^{\frac{b}{a}\left( a- \frac{\alpha}{2} + \frac{a}{b}r \right)n_t } } { (n_t!)^{\frac{b}{a}\left( a - \frac{\alpha}{2} + \frac{a}{b}r \right) } }\right] - c \epsilon
		\\&=  c \log(k \sqrt{\theta}R) - c\frac{b}{a}\left( a - \frac{\alpha}{2} + \frac{a}{b}r \right) \log(c) + c\frac{b}{a}\left( a - \frac{\alpha}{2} + \frac{a}{b}r \right) \liminf_{t \to \infty} \frac{1}{n_t}\log\left(\frac{n_t^{n_t}}{(n_t)!}\right) - c\epsilon
		\\& =   c \log(k \sqrt{\theta}R) - c\frac{b}{a}\left( a - \frac{\alpha}{2} + \frac{a}{b}r \right) \log(c) + c\frac{b}{a}\left( a - \frac{\alpha}{2} + \frac{a}{b}r \right)  - c\epsilon
	\end{align*}
and letting $\epsilon$ tend to $0$ we wee that
	\[
		I(n_t) \ge c \left[ \log(k \sqrt{\theta}R) - \frac{b}{a}\left( a - \frac{\alpha}{2} + \frac{a}{b}r \right) \log(c) + \frac{b}{a}\left( a - \frac{\alpha}{2} + \frac{a}{b}r \right) \right] =: h(c).
	\]
We can easily find the maximal $c$ by solving $h'(c)=0$ and doing so we get $c=(k \sqrt{\theta}R)^{\frac{a}{b\left( a - \frac{\alpha}{2} + \frac{a}{b}r \right)}}$. After plugging in this value for $c$ and replacing $R$ we arrive at the following inequality
	\[
		I(n_t) \ge (k\sqrt{\theta})^{\frac{a}{b\left( a - \frac{\alpha}{2} + \frac{a}{b}r \right)} }(\rho_{new}^{1/2})^{\frac{a}{b\left( a - \frac{\alpha}{2} + \frac{a}{b}r \right)}}
	\]
and simplifying exponents gives us \ref{69ineq}.

$\hfill \qed$

We now can prove Theorem \ref{lbound}.

\noindent \textbf{Proof of Theorem \ref{lbound}:} By proposition \ref{Wnprop2}, for and $p,q >0$ with $p^{-1}+q^{-1} = 1$ we have that
	\[
		\Norm{u(t,0)}_p \ge \exp\left\{ -\frac{1}{2} t_p^\beta \Norm{f}_{\mathcal{H}}^2 \right\}\left| \sum_{n \ge 0} \theta^{n/2} W_n(t_p^\beta, \mathcal{F}f) \right|.
	\]
We now take the supremum over all $f \in \mathcal{H}_+$ with $\Norm{f}_{\mathcal{H}}=k >0 $. Recall that $W_n(t,\phi) = k^nW_n(t,\phi/k)$ for $\phi \in L^2(\mu)$ and non-negativity of $W_n(t,\cdot)$ on $\mathcal{H}_+$. Let $c>0$.
	\begin{align*}
		\Norm{u(t,0)}_p &\ge 	\exp\left\{ -\frac{1}{2} t_p^\beta k^2 \right\} \sup_{f\in \mathcal{H}_+ \; \Norm{f}_{\mathcal{H}}=k}\sum_{n \ge 0} \theta^{n/2} W_n(t_p^\beta, \mathcal{F}f)
		\\&= \exp\left\{ -\frac{1}{2} t_p^\beta k^2 \right\} \sup_{f\in \mathcal{H}_+ \; \Norm{f}_{\mathcal{H}}=1}\sum_{n \ge 0} k^n \theta^{n/2} W_n(t_p^\beta, \mathcal{F}f)
		\\& \ge \exp\left\{ -\frac{1}{2} t_p^\beta k^2 \right\} \sup_{f\in \mathcal{H}_+ \; \Norm{f}_{\mathcal{H}}=1} k^{n_t} \theta^{{n_t}/2} U_{n_t}(t_p^\beta, f)
	\end{align*}
where $n_t = [ct_p^\beta]$. Now by choosing $c$ as in Lemma \ref{logUnlb} we get that
	\begin{align*}
		\liminf_{n \to \infty} t_p^{-\beta} \log \Norm{u(t,x)}_p &\ge 	\liminf_{n \to \infty} \left( \frac{-\frac{1}{2}t_p^\beta k^2}{t_p^\beta} + \frac{1}{t_p^\beta}\log \left[ \sup_{f\in \mathcal{H}_+ \; \Norm{f}_{\mathcal{H}}=1} k^{n_t} \theta^{{n_t}/2} U_{n_t}(t_p^\beta, f) \right]  \right)
		\\& -\frac{1}{2}k^2 + (k\sqrt{\theta})^{ \frac{2a}{2ab-\alpha b +2ar} }(\rho^{1/2})^{\frac{2a}{2ab -\alpha b + 2ar}} =: h(k)
	\end{align*}
By maximizing $h$ we see that $h$ is maximized at the point $k^* = \left( \frac{2ab -\alpha b + 2ar}{2aB} \right)^{ \frac{2ab-\alpha b +2ar}{2a-2[2ab-\alpha b + 2ar]} }$ where $B=(\sqrt{\theta}\rho^{1/2})^{ \frac{2a}{2ab-\alpha b + 2ar} }$. Inserting $k*$ into $h$ gives us that
	\[
		h(k^*) = B^{\beta}\left( \frac{2a}{2ab - \alpha b +2ar} \right)^\beta \left( \frac{2ab-\alpha b + 2ar -a }{2a} \right)
	\]
and plugging in the value for $B$ yields \ref{lbound}.

$\hfill \qed$

\addcontentsline{toc}{section}{Bibliography}
\begin{thebibliography}{999}
		\bibitem{BalanChenChen} % {{{ 1.
		Balan, R. M., Chen, L. and Chen, X. (2021).
		\newblock Exact asymptotics of the stochastic wave equation with time-independent noise.
		\newblock {\em Ann. Inst. Henri Poincar\'e: Probab. \& Stat.}, to appear.
		% }}}
		\bibitem{BalanSong} % {{{ 2.
		Balan, R. M. and Song, J. (2017).
		\newblock Hyperbolic Anderson Model with space-time homogeneous Gaussian noise.
		\newblock {\em ALEA, Latin Amer. J. Probab. Math. Stat.} {\bf 14}, 799--849.
		% }}}

		\bibitem{ChenHuNualartFraction}
		Chen, L., Hu, Y. and Nualart D. (2019).
		\newblock Nonlinear stochastic time-fractional slow and fast diffusion equations on $\R^d$.
		\newblock {\em Science Direct, Stochastic Processes and Their Applications.} \textbf{129}, 5073-5112
		
		\bibitem{Chen2007} Chen X.
		\newblock{\it Parabolic Anderson model with rough or critical Gaussian noise.}
		\newblock{(2019) Ann. Inst. Henri Poincare: Prob. Stat.} \textbf{55}, 941--976.

		\bibitem{chenlarge} Bass, R., Chen, X. and Rosen, M. (2009). Large deviations for Riesz potentials of additive processes. Ann. Inst. Henri Poincare: Probab. $\&$ Stat. 45, 626-666

		\bibitem{ChenHuHuHuang2007} Chen, L., Hu, Y., Hu, G. and Huang, J.
		\newblock{\it Stochastic time-fractional diffusion equations on $\R^d$}
		\newblock Stochastics 89 (1) (2007) 171-206

		\bibitem{DalangEtc09Minicourse} Dalang, Robert, Davar Khoshnevisan, Carl
		Mueller, David Nualart, and Yimin Xiao
		\newblock {\it A minicourse on stochastic partial differential equations}.
		\newblock Held at the University of Utah, Salt Lake City, UT, May 8–19, 2006.
		Edited by  Khoshnevisan and Firas Rassoul-Agha. Lecture Notes in Mathematics,
		1962.
		\newblock Springer-Verlag, Berlin, 2009. xii+216 pp.
		
		\bibitem{HuFracBrown} Hu, Y.
		 Hu, Y. (2001)
		\newblock Heat equation with fractional white noise potential.
		\newblock {\em Appl. Math. Optim.} \textbf{43}, 221-243.
		
		\bibitem{LiebAnalysis} % {{{ 2.
		Lieb, E. and Loss, M. (2001).
		\newblock Analysis
		\newblock {\em Graduate Studies in Mathematics, AMB}.
		% }}}
		
		\bibitem{nualartmalliavin} 
		Nualart D. and Nualart E. (2018)
		\newblock Introduction to Malliavin Calculus.
		\newblock Cambridge University Press. 
		
\end{thebibliography}
\end{document}
